% This must be in the first 5 lines to tell arXiv to use pdfLaTeX, which is strongly recommended.
\pdfoutput=1
% In particular, the hyperref package requires pdfLaTeX in order to break URLs across lines.

\documentclass[11pt]{article}

% Change "review" to "final" to generate the final (sometimes called camera-ready) version.
% Change to "preprint" to generate a non-anonymous version with page numbers.
\usepackage[final]{acl}

% Standard package includes
\usepackage{times}
\usepackage{latexsym}

% For proper rendering and hyphenation of words containing Latin characters (including in bib files)
\usepackage[T1]{fontenc}
% For Vietnamese characters
% \usepackage[T5]{fontenc}
% See https://www.latex-project.org/help/documentation/encguide.pdf for other character sets

% This assumes your files are encoded as UTF8
\usepackage[utf8]{inputenc}

% This is not strictly necessary, and may be commented out,
% but it will improve the layout of the manuscript,
% and will typically save some space.
\usepackage{microtype}

% This is also not strictly necessary, and may be commented out.
% However, it will improve the aesthetics of text in
% the typewriter font.
\usepackage{inconsolata}

%Including images in your LaTeX document requires adding
%additional package(s)
\usepackage{graphicx}

% If the title and author information does not fit in the area allocated, uncomment the following
%
%\setlength\titlebox{<dim>}
%
% and set <dim> to something 5cm or larger.

\usepackage{mdframed}
\usepackage{listings}
\usepackage{xcolor}
\usepackage{array}
\usepackage{booktabs}
\usepackage{multirow}
\usepackage{amsmath}
\usepackage{colortbl}
\usepackage{pifont}
\usepackage{makecell}

\usepackage[ruled,vlined]{algorithm2e}

\definecolor{codebg}{rgb}{0.99,0.99,0.99}
\definecolor{highlightyellow}{rgb}{1,1,0.8}
\definecolor{highlightpink}{rgb}{1,0.8,0.8}
\definecolor{highlightgreen}{rgb}{0.2,0.85,0.2}
\definecolor{darkblue}{rgb}{0.0, 0.0, 0.55}

\definecolor{keyword}{RGB}{0, 0, 255}
\definecolor{comment}{RGB}{0, 128, 0}
\definecolor{string}{RGB}{163, 21, 21}
\definecolor{background}{RGB}{255, 255, 255}
\definecolor{identifier}{RGB}{128, 0, 128}


\newcommand{\lstcolourline}[1]{\rlap{\color{#1}\rule[-.3\baselineskip]{7.35cm}{\baselineskip}}}
\definecolor{mycolor}{rgb}{0.95,0.95,0.95}


\newmdenv[
  backgroundcolor=mycolor,
  linecolor=black,
  linewidth=1pt,
  roundcorner=10pt,
  innerleftmargin=10pt,
  innerrightmargin=10pt,
  innertopmargin=10pt,
  innerbottommargin=10pt
]{custombox}


\lstdefinestyle{mystyle}{
    language=Python,
    basicstyle=\ttfamily\small,
    keywordstyle=\color{blue},
    stringstyle=\color{highlightgreen},
    commentstyle=\color{purple},
    backgroundcolor=\color{codebg},
    numbers=left,
    numberstyle=\footnotesize\ttfamily\hfill, % 右对齐行号
    numbersep=3pt, % 调整行号与代码的距离
    frame=tb,
    rulecolor=\color{black},
    framerule=0.5mm,
    framesep=2mm,
    xleftmargin=4.5mm,
}


\lstdefinestyle{customlst}{
    frame=topline,
    framexbottommargin=5pt,
    framextopmargin=5pt,
    xleftmargin=5pt,
    xrightmargin=5pt,
    aboveskip=10pt,
    belowskip=10pt,
    language=Python,
    basicstyle=\ttfamily,
    keywordstyle=\color{blue}\bfseries,
    commentstyle=\color{gray},
    stringstyle=\color{red},
    morekeywords={Line Number, Type, Explanation, Dead code},
    breaklines=true,
    tabsize=2,
    numbers=left,
    numberstyle=\tiny\color{gray},
    stepnumber=1,
    numbersep=5pt,
    showspaces=false,
    showstringspaces=false
}

\lstdefinestyle{JavaStyle}{
    backgroundcolor=\color{background},
    commentstyle=\color{comment},
    keywordstyle=\color{keyword},
    stringstyle=\color{string},
    basicstyle=\ttfamily\footnotesize,
    breakatwhitespace=false,
    breaklines=true,
    captionpos=b,
    keepspaces=true,
    numbers=none,
    showspaces=false,
    showstringspaces=false,
    showtabs=false,
    tabsize=2,
    language=Java
}

% Define Python language settings
\lstdefinestyle{PythonStyle}{
    keywordstyle=\color{blue},
    stringstyle=\color{highlightgreen},
    backgroundcolor=\color{background},
    commentstyle=\color{comment},
    basicstyle=\ttfamily\footnotesize,
    breakatwhitespace=false,
    breaklines=true,
    captionpos=b,
    keepspaces=true,
    numbers=none,
    showspaces=false,
    showstringspaces=false,
    showtabs=false,
    tabsize=2,
    language=Python
}


\title{DCE-LLM: Dead Code Elimination with Large Language Models}

% Author information can be set in various styles:
% For several authors from the same institution:
% \author{Author 1 \and ... \and Author n \\
%         Address line \\ ... \\ Address line}
% if the names do not fit well on one line use
%         Author 1 \\ {\bf Author 2} \\ ... \\ {\bf Author n} \\
% For authors from different institutions:
% \author{Author 1 \\ Address line \\  ... \\ Address line
%         \And  ... \And
%         Author n \\ Address line \\ ... \\ Address line}
% To start a separate ``row'' of authors use \AND, as in
% \author{Author 1 \\ Address line \\  ... \\ Address line
%         \AND
%         Author 2 \\ Address line \\ ... \\ Address line \And
%         Author 3 \\ Address line \\ ... \\ Address line}

% \author{First Author \\
%   Affiliation / Address line 1 \\
%   Affiliation / Address line 2 \\
%   Affiliation / Address line 3 \\
%   \texttt{email@domain} \\\And
%   Second Author \\
%   Affiliation / Address line 1 \\
%   Affiliation / Address line 2 \\
%   Affiliation / Address line 3 \\
%   \texttt{email@domain} \\}

%\author{
%  \textbf{First Author\textsuperscript{1}},
%  \textbf{Second Author\textsuperscript{1,2}},
%  \textbf{Third T. Author\textsuperscript{1}},
%  \textbf{Fourth Author\textsuperscript{1}},
%\\
%  \textbf{Fifth Author\textsuperscript{1,2}},
%  \textbf{Sixth Author\textsuperscript{1}},
%  \textbf{Seventh Author\textsuperscript{1}},
%  \textbf{Eighth Author \textsuperscript{1,2,3,4}},
%\\
%  \textbf{Ninth Author\textsuperscript{1}},
%  \textbf{Tenth Author\textsuperscript{1}},
%  \textbf{Eleventh E. Author\textsuperscript{1,2,3,4,5}},
%  \textbf{Twelfth Author\textsuperscript{1}},
%\\
%  \textbf{Thirteenth Author\textsuperscript{3}},
%  \textbf{Fourteenth F. Author\textsuperscript{2,4}},
%  \textbf{Fifteenth Author\textsuperscript{1}},
%  \textbf{Sixteenth Author\textsuperscript{1}},
%\\
%  \textbf{Seventeenth S. Author\textsuperscript{4,5}},
%  \textbf{Eighteenth Author\textsuperscript{3,4}},
%  \textbf{Nineteenth N. Author\textsuperscript{2,5}},
%  \textbf{Twentieth Author\textsuperscript{1}}
%\\
%\\
%  \textsuperscript{1}Affiliation 1,
%  \textsuperscript{2}Affiliation 2,
%  \textsuperscript{3}Affiliation 3,
%  \textsuperscript{4}Affiliation 4,
%  \textsuperscript{5}Affiliation 5
%\\
%  \small{
%    \textbf{Correspondence:} \href{mailto:email@domain}{email@domain}
%  }
%}

\author{
 Minyu~Chen$^{1}$, Guoqiang Li$^{1}\thanks{~~Corresponding author.}$ , Ling-I Wu$^{1}$, Ruibang Liu$^{1}$ \\
 $^{1}$Shanghai Jiao Tong University, Shanghai, China\\
 \texttt{{\{minkow,li.g,edithwuly,628628\}@sjtu.edu.cn}}
}

\begin{document}
\maketitle
\begin{abstract}
Dead code introduces several challenges in software development, such as increased binary size and maintenance difficulties. It can also obscure logical errors and be exploited for obfuscation in malware. For LLM-based code-related tasks, dead code introduces vulnerabilities that can mislead these models, raising security concerns.
Although modern compilers and IDEs offer dead code elimination, sophisticated patterns can bypass these tools. A universal approach that includes classification, location, explanation, and correction is needed, yet current tools often require significant manual effort.
We present DCE-LLM, a framework for automated dead code elimination using a small CodeBERT model with an attribution-based line selector to efficiently locate suspect code. LLMs then generate judgments and explanations, fine-tuned on a large-scale, annotated dead code dataset to provide detailed explanations and patches. DCE-LLM outperforms existing tools, with advanced unreachability detection, automated correction, and support for multiple programming languages. Experimental results show DCE-LLM achieves over 94\% F1 scores for unused and unreachable code, significantly surpassing GPT-4o by 30\%.\footnote{Our dataset and code are available at \url{ https://github.com/Minkow/DCE-LLM}}
\end{abstract}

\section{Introduction}
% Dead code, defined as segments of a program that are either never executed or do not affect the program’s overall functionality, poses significant challenges in software development~\cite{8370748}. It contributes to increased binary size, which can degrade performance through heightened instruction cache misses. In languages like JavaScript~\cite{malavolta2023javascript}, although dead code does not execute at runtime, it still incurs performance costs during the downloading and parsing phases, further prolonging execution times. Additionally, dead code complicates software maintenance by cluttering the codebase, thus consuming developers' time as they attempt to comprehend its purpose. More critically, logical errors can hide behind unreachable statements caused by flawed conditional logic or incorrect assumptions in the code flow, inadvertently masking the presence of bugs. Malware authors also leverage dead code as an obfuscation technique to increase the overhead of reverse analysis.

Dead code, defined as segments of a program that are never executed or that do not affect program functionality, can introduce performance bottlenecks and obscure potential security vulnerabilities. The presence of dead code increases binary size, which negatively impacts system performance, particularly in resource-constrained environments like web applications~\cite{8370748,malavolta2023javascript}. Additionally, dead code complicates software maintenance by cluttering the codebase, thus consuming developers' time as they attempt to comprehend its purpose. In more critical cases, dead code can serve as a vector for obfuscation techniques employed in adversarial attacks, complicating code comprehension and analysis.

Though large language models (LLMs) have achieved impressive results in various code-related tasks~\cite{weicoeditor, jainllm, Tang2024CodeAgentAC, li2024rewriting, han2024archcode}, they are particularly vulnerable when faced with dead code. Dead code, whether introduced unintentionally or as part of an attack, can disrupt a model's ability to accurately interpret program logic. In particular, dead code injection attacks exploit this vulnerability by adding redundant or unreachable code, which confuses the model and leads to incorrect predictions or faulty analyses. Studies have shown that this can cause a significant drop in accuracy for tasks like vulnerability detection, with an observed reduction of up to 12.7\%~\cite{khare2023understanding}. These weaknesses highlight the need for more robust approaches to handling dead code in LLM-based tasks.
% The dead code injection attacks can exploit dead code to mislead these models, resulting in undesired output or prediction. A recent study points out that LLMs suffer performance degradation in the code vulnerability detection task while facing dead code insertion attacks, showing an average accuracy reduction of up to 12.7\%~\cite{khare2023understanding}. This not only jeopardizes the models' accuracy but also raises security concerns in future software engineering applications based on language models.

% While modern compilers can effectively remove dead variables from programs~\cite{wang2017dead, romano2016graph}, and certain IDEs can alert programmers to the presence of dead code during the compiling stage, the dead code elimination process remains fraught with challenges. Firstly, sophisticated attack patterns based on unreachable statements can bypass these tools, potentially leading to severe security vulnerabilities and performance issues. These well-designed attacks can embed dead code that remains undetected, exploiting the weaknesses of language models and compromising the integrity of software systems. Moreover, the phenomenon of dead code is pervasive across all programming languages, underscoring the need for a universal dead code elimination method. Each programming language has its own syntax and semantics, which complicates the development of a single solution that can effectively handle dead code in every context. Thus, a versatile approach is essential for the robustness and maintainability of software engineering. Nevertheless, a complete dead code elimination framework should encompass the processes of classification, location, explanation, and correction. However, the explanation and correction stages are particularly challenging for current program analysis tools, often requiring considerable human effort to complete. An automated method capable of executing the entire dead code elimination framework is highly desirable.

While modern compilers can remove basic dead variables and IDEs can flag some dead code~\cite{wang2017dead, romano2016graph}, effectively eliminating all forms of dead code remains a challenge. Moreover, sophisticated attack patterns using unreachable statements can bypass these tools, leading to potential security vulnerabilities and performance issues. These attacks exploit weaknesses in language models, embedding dead code that remains undetected. Additionally, the pervasiveness of dead code across different programming languages complicates the development of a universal solution, as each language has its own syntax and semantics. As a result, a robust and versatile dead code elimination framework is essential, one that can automate the processes of classification, detection, explanation, and correction, thereby reducing the need for manual intervention.
% Furthermore, a complete dead code elimination framework must encompass classification, location, explanation, and correction. However, the explanation and correction stages are particularly challenging for current program analysis tools, introducing considerable human effort. Therefore, automating these stages is highly desirable to enhance efficiency and streamline the entire elimination process.

% Despite advancements in modern compilers and integrated development environments (IDEs) that can eliminate straightforward cases of dead code, complex patterns, especially those introduced via unreachable conditions or adversarial dead code injection, remain difficult to detect. These patterns are designed to evade traditional static analysis tools, necessitating more advanced, automated methods for comprehensive dead code elimination.

% While modern Compilers can effectively remove dead variables from programs, and certain IDEs can alert programmers to the presence of dead code during the compiling stage, the dead code elimination process remains fraught with challenges. Firstly, sophisticated attack patterns based on unreachable statements can bypass these tools, potentially leading to severe security vulnerabilities and performance issues. These well-designed attacks can embed dead code that remains undetected, thus exploiting the weaknesses of language models and compromising the integrity of software systems. Moreover, the phenomenon of dead code is pervasive across all programming languages, underscoring the need for a universal dead code elimination method. Each programming language has its own syntax and semantics, which complicates the development of a single solution that can effectively handle dead code in every context. Thus, a versatile approach is essential for the robustness and maintainability of software engineering. Nevertheless, the complete dead code elimination framework should contain location-explaination-correction processes. However, the explanation and correction processes are challenging for current program analysis tools. Considerable human efforts are paid to accomplish the correction stage. Thus, a () method are desired to automatically execute the complete dead code elimination framework.

% Although LLMs can be fooled by dead code snippets, they have shown stunning performance on various natural language and programming language tasks. Those LLMs can even perform human-like program analysis processes such as (). We believe that language models pre-trained on the code corpus have obtained essentials for addressing dead code issues. For unused codes including unused variables and irrelevant operations, language models can (), for unreachable statements, language models can (). However, LLMs are not well-suited for dead code elimination for several reasons. First, LLMs are trained on publicly available programs with dead codes. In our study, 17.2\% python code in the CodeSearchNet dataset contains dead code, but no specialized indicates for dead code are given. Second, LLMs may ignore dead code.
% LLMs struggle to robustly access and use information split in long input contexts, especially when handling dead code that can occur in every line. Thirdly, it is challenging for LLMs to explain and correct dead codes. (reasons)

% Although LLMs can be fooled by dead code snippets, they have shown impressive performance on various natural language and programming language tasks. These models can even perform human-like program analysis processes such as code completion, refactoring, and debugging. We believe that LLMs pre-trained on large code corpora have acquired the essentials for addressing dead code issues. For example, LLMs can identify unused code, such as variables and irrelevant operations, because they learn patterns and dependencies in code during training, enabling them to recognize code segments that do not contribute to the program’s functionality. For unreachable statements, LLMs are able to analyze conditional expressions and evaluate their logical outcomes, similar to how symbolic execution tools operate. However, LLMs are not well-suited for dead code elimination for several reasons. First, LLMs are trained on publicly available programs that often contain dead code. For example, over 45.3\% of Java code in the CodeNet~\cite{puri2021codenet} dataset contains dead code, but there are no specialized processing steps for filtering dead code before training. Second, LLMs may ignore dead code because they struggle to robustly access and use information split across long input contexts, especially when handling dead code that can occur in every line. Thirdly, explaining and correcting dead code is challenging for LLMs, as they require a deep understanding of the program's intended functionality, context, and the ability to generate semantically correct patches.




% Consequently, an integrated approach that combines the strengths of LLMs with advanced static and dynamic analysis techniques is necessary to effectively tackle dead code elimination.

% because an unignored portion of publicly available code in the training corpus contains dead code, but no specialized indicates for dead code are given. To unlock the LM's capability for dead code elimination,


% Considering the limitations of existing tools and the challenges associated with applying LLMs to dead code elimination, we make a significant step towards addressing these issues with DCE-LLM, an LLM-empowered framework for comprehensive dead code elimination. To tackle the above-mentioned issues while directly adopting LLMs for dead code elimination, we first leverage a relatively small CodeBert model with a novel attribution-value-based line selector to effectively and accurately locate suspect dead code snippets in long code inputs.
% This step filters out most normal code to reduce the call time of LLMs. Then, LLMs are employed to generate the final judgment along with explanations of dead code given the source code and highlighted lines of suspects. This allows LLMs to focus on a selected part of the program. Furthermore, we utilize fine-tuning techniques on LLMs with the first large-scale dead code dataset which is automated annotated, equipping LLMs with the capability of generating explanations and patches for dead code. As the first neural-based framework for dead code elimination, DCE-LLM demonstrates several advantages compared to existing tools:

% \textit{Unreachability Checking. } DCE-LLM is capable of detecting complex unreachable dead code. This type of dead code can be employed to conduct dead code insertion adversarial attacks, bypassing most compilers and linters.

% \textit{Explanation and Correction. } DCE-LLM can automatically provide explanations and patches for dead code, beyond simply locating it. This significantly reduces the manual effort required by developers, improving the overall efficiency of the development process.

% \textit{Language Support. } DCE-LLM benefits from the natural ability of LLMs to understand and generate code in multiple programming languages. DCE-LLM can perform dead code elimination on diverse programming languages including Python and Java, thereby enhancing its versatility and applicability across different development environments. Besides, DCE-LLM also obtains the capability of analyzing incomplete code such as standalone methods without building projects, while program analysis-based approaches failed to construct control-flow graphs (CFGs) or data-flow graphs (DFGs) due to missing declared data types or missing library imports.


% Our experimental results show that DCE-LLM achieved recall rates over 95\% for both unused and unreachable code, along with over xxx\% F1 rates, which significantly outperform LLMs and IDEs. Furthermore, in terms of the quality of generated content including explanation and repaired code, DCE-LLM successfully beat the baseline models in an arena-like human evaluation.

% We highlight our contribution as follows:
% \begin{enumerate}
%     \item \textbf{Significance: }We propose a novel framework, DCE-LLM, for employing LLMs to address the dead code elimination task including classification, location, exploration, and patching. To the best of our knowledge, we are the first to employ LLMs for dead code elimination, which offers complex unreachability checking, explanation, and code repairing.
%     \item \textbf{Novelty: }We employ a relatively small model along with a novel attribution technique to effectively locate suspect lines. LLMs are empowered with valuable information provided by the pivot model.
%     \item \textbf{Usefulness: }We propose the first large-scale dataset aimed at dead code elimination. With this dataset, we successfully demonstrate that DCE-LLM can almost perfectly detect unused and unreachable code in multiple programming languages with over 95\% F1 score.
% \end{enumerate}

% An Organization Paragraph can be deleted.

LLMs face several key challenges when it comes to dead code elimination. First, LLMs are trained on publicly available code datasets, many of which contain dead code. However, these datasets lack specific labels for dead code, meaning LLMs do not explicitly learn to identify or handle it during training. For example, we find that over 45.3\% of Java code in the CodeNet~\cite{puri2021codenet} dataset contains dead code, but this is not filtered or annotated, leading to gaps in the models' understanding.
Second, LLMs struggle to locate dead code in long inputs. Dead code can appear in various parts of a program, often scattered across lines of code. LLMs have difficulty maintaining robust context over long inputs, making it challenging to accurately detect and address dead code in complex or extended codebases.
Third, while LLMs excel at tasks like code generation and refactoring, explaining and correcting dead code presents a different challenge. Properly addressing dead code requires a deep understanding of the intended functionality and context of the program, which LLMs may fail to grasp. Additionally, generating semantically correct patches for dead code requires advanced reasoning capabilities that go beyond simple code completion or modification tasks.



To address these challenges, we present DCE-LLM, an LLM-empowered framework for comprehensive dead code elimination. We leverage a relatively small CodeBERT model with a novel attribution-value-based line selector to effectively and accurately locate suspect dead code snippets in long code inputs. This step filters out most normal code, reducing the call time for LLMs. Subsequently, LLMs generate final judgments and explanations for dead code, focusing on the highlighted suspect lines. Furthermore, we fine-tune LLMs on the first large-scale dead code dataset, which is automatically annotated, equipping them with the capability to generate explanations and patches for dead code. As the first neural-based framework for dead code elimination, DCE-LLM demonstrates several advantages over existing tools:

\textit{Unreachability Checking.} DCE-LLM effectively detects complex unreachable dead code that can bypass compilers and IDEs, addressing vulnerabilities used in adversarial attacks.

\textit{Explanation and Correction.} DCE-LLM automatically provides explanations and patches for dead code, significantly reducing the manual effort required by developers and enhancing overall development efficiency.

\textit{Programming Language Support.} DCE-LLM leverages the natural ability of LLMs to understand and generate code in multiple programming languages. It supports multiple programming languages, including Python and Java, making it applicable across diverse development environments.

Our experimental results demonstrate that DCE-LLM achieves F1 scores over 94\% for both unused and unreachable code, along with F1 scores significantly surpassing those of existing LLMs and IDEs. Moreover, in terms of the quality of generated content, including explanations and repaired code, DCE-LLM outperformed baseline models in arena-style human evaluations.

We highlight our contributions as follows:
\begin{itemize}
\item We propose a novel framework, DCE-LLM, for leveraging LLMs to tackle the dead code elimination task, encompassing classification, location, exploration, and patching. To our knowledge, this is the first application of LLMs for dead code elimination, offering sophisticated unreachability checking, detailed explanations, and effective code repair.
\item We employ a small CodeBERT model with an innovative attribution technique to accurately locate suspect lines, augmenting LLMs with extra those hints.
\item We introduce the first large-scale dead code dataset, demonstrating that DCE-LLM achieves over 94\% F1 scores in detecting unused and unreachable code across multiple programming languages, surpassing GPT-4o with over 30\% F1.
\end{itemize}


\section{Related Work}
\textbf{Dead Code Elimination}. Researchers have proposed various approaches for detecting dead codes. Chen et al.~\cite{chen1998c++} introduced a data model serving reachability
analysis and dead code detection for C++ repositories. Boomsma et al.~\cite{boomsma2012dead} leverages a dynamic framework for extracting dead files in PHP web apps by monitoring file usage. Several works focus on call graph analysis~\cite{romano2016graph, romano2018exploring} while program slicing is also adopted to build generic frameworks~\cite{alabwaini2018using, wang2017dead}
. Very recently, Lacuna~\cite{malavolta2023javascript} has presented to integrate third-party analysis techniques for JavaScript dead code detection.

Our proposed DCE-LLM takes a fundamentally different approach.  As a learning-based method, it avoids reliance on specific programming languages or external code analysis tools like LLVM~\cite{wang2017dead}, or hand-crafted rules. This allows DCE-LLM to generalize across multiple languages and handle incomplete or even non-compilable code, making it both versatile and robust.  Critically, our method effectively addresses the challenge of adversarial unreachable code and offers practical solutions for dead code repair, completing the dead code elimination pipeline rather than just performing detection.

% \textbf{LLMs for Code.} Besides researches focusing on code generation and code completion, recent advances in LLMs have shown impressive results, particularly in code-related tasks such as bug detection~\cite{wang-etal-2024-sanitizing}, code transpilation~\cite{bhatia2024verified} and code repair~\cite{tang2024code, zhao-etal-2024-repair}. NS-Slicer~\cite{yadavally2024learning} fine-tuned a pre-trained LLM to predict static program slices for both complete and partial code and leveraged its predictions for vulnerability detection. These works demonstrate the code reasoning capabilities of LLMs across various scenarios, forming an important foundation for the effectiveness of DCE-LLM.


\textbf{LLMs for Code.} LLMs have demonstrated remarkable capabilities in various code-related tasks beyond generation and completion. Recent research highlights their effectiveness in areas like bug detection~\cite{wang-etal-2024-sanitizing}, code transpilation~\cite{bhatia2024verified}, and code repair~\cite{tang2024code, zhao-etal-2024-repair}. For example, NS-Slicer~\cite{yadavally2024learning} fine-tuned an LLM to predict static program slices for both complete and partial code, subsequently using these predictions to enhance vulnerability detection. These advancements showcase the code reasoning abilities of LLMs in diverse scenarios and provide a strong foundation for the effectiveness of DCE-LLM.

\section{Background}

% \subsection{Dead Code Elimination}
\begin{figure}[!bp]
\begin{minipage}{\linewidth}
\small
\begin{lstlisting}[style=mystyle, escapechar=!]
def fill_str(Data):
  s1 = input()
  s2 = s1 + '<PAD>'
!\lstcolourline{yellow!25}!  s3 = s1 + '<EOS>'  # Unused Variable
!\lstcolourline{pink!20}!  if len(s2) == 0: # Unreachable Code
    print('Empty string')
    Data.pad_str = None
    Data.eos_str = None
  else:
    Data.pad_str = s2
    Data.eos_str = 's3'
\end{lstlisting}
\end{minipage}
\caption{An illustrative Python method with dead code.}
\label{fig1}
\end{figure}


As illustrated in Figure \ref{fig1}, Python method \texttt{fill\_str} contains several issues about two primary types of dead code. \textbf{Unused code} refers to code defined or executed but whose result is never used in any other computation. The execution of dead code wastes computation time and memory. In our example, The author of this method has accidentally put quotes around \texttt{s3} in line 11, resulting in an unused variable \texttt{s3} in the fourth line in light yellow.
\textbf{Unreachable code}, defined as code that can never be executed. Line 5 in light pink is the start of an unreachable code snippet. It introduces an unreachable case caused by an always-false condition. Thus, code in lines 6-8 cannot be executed whatever user input in line 2.

Modern compilers and IDEs offer extensive support for dead code elimination, especially for unused code. However, detecting unreachable code is inherently challenging, especially those designed for obfuscation or adversarial attacks as illustrated in Figure \ref{fig1}.
% In our example, the condition \texttt{if len(s2) == 0:} is always \texttt{false} because \texttt{s2} is initialized as \texttt{s1 + `<PAD>'}.
Static analysis tools might not catch this unreachable code \texttt{if len(s2) == 0:} if the condition requires to be evaluated with a satisfiability check. Moreover, several attack patterns are introduced~\cite{gao2023discrete} for producing complex unreachable code branches which are imperceptible to compilers and IDEs. They successfully mislead the output result of language models such as CodeBERT. Even powerful LLMs like GPT-4o are occasionally affected, without realizing that the dead branch is unreachable in the vulnerability detection task~\cite{khare2023understanding}.
% Besides, dead code eliminators are usually language-specific and integrated into different compilers and IDEs, which require extensive human efforts for designing and testing dead code patterns across multiple programming languages.

% \subsection{Large Language Models}

% Large language models (LLMs) are extensively employed in software engineering, fulfilling two primary roles. Firstly, many researchers exploit the remarkable code generation capabilities of these models, treating them as generators. LLMs are used to convert natural language descriptions of software into executable code~\cite{rao2023cat, he2024conline}, generate test cases~\cite{deng2024large}, and enhance existing tools like SMT solvers~\cite{sun2023smt}. However, it has been noted that uncontrolled bugs may be introduced during the generation and translation process~\cite{pan2024lost}. Secondly, due to their strong code understanding abilities, LLMs are also employed as detectors or analyzers. This includes applications in fault and vulnerability localization~\cite{qin2024agentfl, zhang2024empirical, yang2024large, khare2023understanding} and program slicing~\cite{yadavally2024learning}. Furthermore, LLMs are utilized in complex tasks that combine both aspects, such as program repair~\cite{huang2023empirical, hossain2024deep}, showcasing their versatility and depth in enhancing software quality and reliability.

% As the capabilities of large language models in code generation and code understanding tasks have been demonstrated, it is time to move on to a more complex task: dead code elimination, which is a critical topic in software engineering. Boomsma et al.~\cite{boomsma2012dead} have proposed an identification and removal approach based on dynamic analysis for web systems written in PHP. Theodoridis et al.~\cite{theodoridis2022finding} has developed an effective dead code detection approach using control flow graphs to analyze how well compilers optimize code and identify missed optimizations, showing that even widely-used compilers like GCC and LLVM have room for improvement in dead code optimization. Malavolta et al.~\cite{malavolta2023javascript} have introduced an approach for automatically detecting and eliminating JavaScript dead code from web apps, supporting both static and dynamic analyses. For the dead code elimination task, code understanding capabilities are required to comprehend the execution logic and recognize dead code. Simultaneously, code generation capabilities are necessary to produce clean code without altering the execution logic as well.

% Large language models (LLMs) are ... This part for introducing LLMs in SE.

% Hence, with text-based prompts containing source codes and task specifications, LLMs can identify some cases of dead code. Figure \ref{fig2} shows the result of requesting GPT-4o to locate the dead code snippets presented in Figure \ref{fig1}. This prompt can be regarded as a baseline of dead code elimination with LLMs.

% \begin{figure}[ht]
% \begin{minipage}{\linewidth}
% \small
% \noindent\makebox[\linewidth]{\rule{\linewidth}{1pt}}
% \colorbox{codebg}{
% \parbox{\dimexpr\linewidth-2\fboxsep}{
% \textbf{\textcolor{cyan}{Input:}} \\
% You are an experienced programmer expert in detecting dead code such as unused code and unreachable code. Given the following Python code, respond in the following format: \\
% \textcolor{orange}{Dead code:} Yes or No \\
% For each dead code,\\
% \textcolor{gray}{Line Number:} Line number\\
% \textcolor{olive}{Type:} Unused or Unreachable\\
% \textcolor{purple}{Explanation:} Explanation for prediction\\

% Finally, fix the code with: \\
% \textcolor{brown}{Fixed Code:} Fixed code \\

% \textcolor{blue}{Code:} [Code in Figure 1] \\

% \textbf{\textcolor{magenta}{Output:}} \\
% \textcolor{orange}{Dead code:} \textcolor{red}{Yes} \\
% \textcolor{gray}{Line Number:} \textcolor{red}{4} \\
% \textcolor{olive}{Type:} Unused \\
% \textcolor{purple}{Explanation:} The variable s3 is defined but never used in any subsequent code. Instead of using s3, the code mistakenly uses the literal string `s3'. \\

% \textcolor{brown}{Fixed Code:} [Code in Figure 1 after fix line 11]
% }}
% \noindent\makebox[\linewidth]{\rule{\linewidth}{1.5pt}}



% \end{minipage}
% \caption{The basic prompt and a sample response for DCE with GPT-4o.}
% \label{fig2}
% \end{figure}

% In this prompt example, we provide GPT-4o, one of the most powerful LLMs, with the code snippet to be analyzed. We expect the LLM to (1) locate where the dead code occurs. (2) Judge the type of the dead code. (3) Explain the reason that this line of code is a dead one. (4) Generate a proper patch to remove the dead code. In this case,
% GPT-4o can understand the request to detect dead code and points out that Line 4 contains unused code snippets. It even provides a possible root cause of this line of dead code. While this dead code occurs in line 4, GPT-4o can generate a proper code patch by fixing line 11. Thus, in our task, GPT-4o shows its promising capability to provide convincing evidence and eliminate dead code.

% However, several issues are raised while directly applying LLMs for DCE.
% First, during some of our tests, the LLMs ignored the human-generated unreachable control branch in Line 5. LLMs tend to be more sensitive to unused rather than unreachable code, leading to performance degradation in code-related tasks as reported in ~\cite{khare2023understanding}.
% Second, LLMs struggle to accurately identify the exact line containing the dead code. The explanations and line numbers often do not align due to the LLMs' weakness in number counting.
% Third, LLMs are prone to generating false positives, mistakenly identifying normal code as unused or unreachable.
% Fourth, since most code snippets are normal and do not contain dead code, processing every snippet with LLMs can be computationally expensive, especially for open-source models like Llama~\cite{touvron2023llama}. As for close-source LLMs such as the GPT family, the API cost can be substantial.

% From this running example, we can see that LLMs still lack the capability to locate all dead code snippets. Another important issue when applying LLMs to DCE is the computational cost for reference with open-source LLMs, or token cost for calling close-source LLMs.

% Given the issues identified, we propose DCE-LLM, an innovative approach that utilizes a relatively small model to accurately extract lines and types of dead code within code snippets. This small model serves as a pivot for LLMs, enhancing their precision and effectiveness in dead code elimination while reducing the false positive rate. Additionally, we harness the generation capabilities of LLMs to provide detailed explanations for dead code and suggest valid patches for its removal.
% To the best of our knowledge, we are the first to present a comprehensive pipeline for identifying, locating, explaining, and eliminating dead code. Moreover, we pioneer the employment of LLMs to manage not only dead code that arises during development but also dead code employed as adversarial attack patterns.

\section{Methodology}

\begin{figure*}[!ht]
\centering
\includegraphics[width=16cm]{dcellm.pdf}
\caption{Overview of the DCE-LLM approach. The yellow labels highlight the training step while the green labels introduce the inference step.}
\label{fig3}
\end{figure*}

% In this section, we introduce our DCE-LLM approach to automatically accomplish dead code elimination. Figure \ref{fig3} illustrates an overview of DCE-LLM. During the training process of DCE-LLM (labeled with yellow backgrounds), we first build a dead code dataset to train a pivot model based on CodeBERT. Then, we request GPT-4o to provide auxiliary
% annotations including explanations and improved code snippets. The annotated dataset is employed to fine-tune an LLM for better dead code detection accuracy and code generation quality. For the inference step (labeled with green backgrounds), we leverage our trained pivot model to efficiently filter out suspect code snippets. Then, we leverage a novel technique, named dead code Shapley, to locate possible dead code lines with effectiveness and efficiency. Finally, we provide detailed prediction results from the pivot model to our fine-tuned LLMs, to generate human-readable reports for dead code and a (有用的) suggestion for dead code elimination.

% As Fig. \ref{fig3} illustrates, we introduce our DCE-LLM approach for automating dead code elimination.
% During the training phase (highlighted in yellow), we start by constructing a dead code dataset to train a pivot model based on CodeBERT. Next, we leverage GPT-4o to provide auxiliary annotations, including detailed explanations and improved code snippets. This annotated dataset is then used to fine-tune an LLM, thereby enhancing both dead code detection accuracy and code generation quality.
% In the inference phase (highlighted in green), our trained pivot model efficiently filters out suspect code snippets. We then apply a novel technique, called dead code attribution, to accurately and efficiently locate potential dead code lines. Finally, the prediction results from the pivot model are leveraged by our fine-tuned LLMs as extra hints to offer human-readable and practical suggestions.
% generate human-readable reports on dead code, offering practical suggestions for its elimination.
This section introduces DCE-LLM (Figure \ref{fig3}), our novel approach for automated dead code elimination, comprising training (yellow) and inference (green) phases.  In training, we first train a CodeBERT-based pivot model for high recall on a constructed dead code dataset.  We then leverage GPT-4o to generate high-quality annotations, including explanations and improved code, for fine-tuning an LLM to enhance detection accuracy and code suggestion quality.  Inference begins with the pivot model filtering suspect code.  Our novel dead code attribution technique then precisely locates and incorporates potential dead code lines into LLM prompts. Finally, the LLM uses pivot model predictions as hints to generate practical dead code removal suggestions.

\subsection{Training Phase}

\subsubsection{Data Collection}
Systematically gathering a comprehensive dataset of both normal and dead code snippets is essential for our approach. This enables the model to learn the intricate details of code structure and functionality, improving its ability to detect and eliminate dead code precisely.

Unluckily, there are no large-scale datasets available specifically for dead code. Thus, we create a novel dataset, AIDCE, aimed at providing a comprehensive collection of annotated both unused and unreachable dead code samples.

In practice, we selected CodeNet~\cite{puri2021codenet} as our data source. CodeNet collects data from online judge websites, consisting of submissions to various programming problems. This dataset offers a rich collection of code snippets in various languages including C++, Python, and Java.
% Compared to code from open-source projects hosted on GitHub, code from online judge platforms tends to be shorter, avoiding the context length limitations of language models. Additionally,
Since no code review process is conducted on online judge platforms, the code submissions often contain more dead code fragments written by programmers. From this corpus, we select code snippets written in Python and Java, resulting in about 300K code files.

% We then leverage existing tools to annotate \textbf{unused} codes. For Python code, we select vulture~\cite{Vulture} to seek dead code out. Vulture is an open-source static code analyzer that can detect unused functions, imports, and variables. Vulture also supports finding unreachable code, but only limited to code after return, and testing simple conditions with \textit{eval()} method in Python. For Java code, we utilize IntelliJ IDEA Java IDE from JetBrains~\cite{intellij2011most} to mark dead codes. IntelliJ IDEA provides a command-line inspector that launches its instance in the background where it runs the inspections. Believing in the power of static analysis, we regard the unused codes detected by the two tools above as gold labels.

We then leverage existing tools to annotate \textbf{unused} code. For Python, we use \textit{Vulture}~\cite{Vulture}, an open-source static code analyzer that detects unused functions, imports, and variables. Vulture can also identify unreachable code, but its capabilities are limited to code following a return statement and simple conditions tested with the \textit{eval()} method in Python. For Java, we utilize \textit{IntelliJ IDEA} from JetBrains~\cite{intellij2011most}. \textit{IntelliJ IDEA} offers a command-line inspector that operates in the background to perform inspections and highlight unused code. By harnessing the power of static analysis, we treat the unused code identified by these tools as gold-standard labels for our dataset.

% However, static analysis-based tools shows vulrunbilty when facing unreachable branch insertion attacks like those presented in Figure \ref{fig1}. These branches or loops whose conditions are always evaluated to false. non-obvious that is complicated enough to bypass compilers and checkers. Besides some attack patterns that fool nearly all compilers and checkers as DaK demonstrates~\cite{gao2023discrete}, we extend the types of patterns to 62. We further demonstrate the effectiveness of those unreachable codes against compilers and checkers in Section \ref{pattern}. We conduct unreachable code insertion for randomly selected files in our code corpus while preserving the original function of those snippets. The inserted dead code blocks are also randomly generated to avoid over-fitting on too specific patterns. In this way, we obtain our \textbf{unreachable} code corpus. In practice, to avoid label leakage and over-fitting, we choose 32 of 62 patterns as training patterns, leaving 30 patterns for only testing.

\begin{table}[tbph]
    \centering
    \renewcommand{\arraystretch}{0.8}
    \small
    \begin{tabular}{>{\raggedright\arraybackslash}p{1cm}lp{0.8cm}p{0.8cm}}
        \toprule
        \multicolumn{1}{c}{\textbf{Name}} & \multicolumn{1}{c}{\textbf{Code}} & \textbf{Python} & \textbf{Java} \\
        \midrule
After return & {\begin{lstlisting}[style=PythonStyle]
return
{deadcode}
\end{lstlisting}}& {\textcolor{brown}{\ding{51}}} {\textcolor{darkblue}{\ding{51}}} & {\textcolor{brown}{\ding{51}}} {\textcolor{darkblue}{\ding{51}}}  \\
\midrule

        Covered branch & {\begin{lstlisting}[style=PythonStyle]
if a>b: {...}
elif a<=b: {...}
else: {deadcode}
\end{lstlisting}}& {\textcolor{brown}{\ding{55}}} {\textcolor{darkblue}{\ding{55}}} & {\textcolor{brown}{\ding{55}}} {\textcolor{darkblue}{\ding{51}}}  \\
\midrule

        Floor & {\begin{lstlisting}[style=PythonStyle]
b = math.floor(a)
if a<b: {deadcode}
\end{lstlisting}}& {\textcolor{brown}{\ding{55}}} {\textcolor{darkblue}{\ding{55}}} & {\textcolor{brown}{\ding{55}}} {\textcolor{darkblue}{\ding{55}}}  \\
\midrule

After assert & {\begin{lstlisting}[style=PythonStyle]
assert a>0
if a<0: {deadcode}
\end{lstlisting}}& {\textcolor{brown}{\ding{55}}} {\textcolor{darkblue}{\ding{55}}} & {\textcolor{brown}{\ding{55}}} {\textcolor{darkblue}{\ding{55}}}  \\
\midrule

Sorted array & {\begin{lstlisting}[style=PythonStyle]
a = sorted([...])
if a[0]>a[-1]:
    {deadcode}

\end{lstlisting}}& {\textcolor{brown}{\ding{55}}} {\textcolor{darkblue}{\ding{55}}} & {\textcolor{brown}{\ding{55}}} {\textcolor{darkblue}{\ding{55}}}  \\


\midrule
\multicolumn{4}{c}{ $\cdots$}\\

\midrule
\midrule
\multicolumn{4}{c}{\textcolor{brown}{Vulture: 1/62} \ \textcolor{darkblue}{PyCharm: 1/62} \ \ \  \textcolor{brown}{JIT: 1/62} \ \textcolor{darkblue}{IDEA: 18/62}}\\


\bottomrule
\end{tabular}
\caption{Examples of unreachable patterns}
\label{taba}
\end{table}

However, static analysis tools demonstrate vulnerabilities when encountering unreachable branch insertion attacks, as shown in Table \ref{taba}. These branches or loops contain conditions that always evaluate to false, yet are often complex enough to evade detection by compilers and analysis tools. In addition to attack patterns that deceive nearly all compilers and checkers, as illustrated by DaK~\cite{gao2023discrete}, we have expanded the range of such patterns to 62, most of which successfully bypass both IDEs and compilers.\footnote{More details can be found in the experienment part.\ref{unreach}}

We perform unreachable code insertion for randomly selected files in our code corpus while preserving the original function of those snippets. The inserted dead code blocks are also randomly generated. To mitigate label leakage and overfitting, we only use 32 out of the 62 patterns for training, reserving the remaining 30 patterns exclusively for testing. Additionally, we keep the original code snippets in the training corpus as hard negatives.


After combining samples of both unused and unreachable code, we construct the AIDCE dataset, which comprises Python and Java code. The AIDCE dataset supports a 3-type classification of dead code and provides auxiliary information such as the specific lines of dead code. The statistics of the AIDCE dataset are presented in \ref{tab:aidce_stats}. The train/test/dev sets are split into approximately 80\%/10\%/10\% respectively.

\begin{table}[h]
\centering
\caption{AIDCE dataset statistics}
\small
\begin{tabular}{lccc}
\toprule
 & Java & Python & \cellcolor{gray!20}Total \\
\midrule
Normal & 4427 & 8201 & \cellcolor{gray!20}12628 \\
Unused & 1853 & 2642 & \cellcolor{gray!20}4495 \\
Unreachable & 207 & 2309 & \cellcolor{gray!20}2516 \\
\cellcolor{gray!20}Total & \cellcolor{gray!20}6487 & \cellcolor{gray!20}13152 & \cellcolor{gray!20}19639 \\
\bottomrule
\end{tabular}
\label{tab:aidce_stats}
\end{table}

% As we expect not just a classifier but a more powerful model supporting locating, explaining, and removing dead code. We look for LLMs to generate human-readable contents supporting dead code elimination with minimal human efforts. Thus, we construct a more detailed and complex training set for LLM fine-tuning. The basic template of dialog is presented in Figure \ref{fig1} and the source of various fields of data in shown in
% \ref{tab1}. Most fields are already filled in the previous step of constructing our AIDCE dataset. Then, we leverage the state-of-the-art GPT-4o model as an annotator to finish the unfilled contents including explanations and fixed code snippets. Some explanations are already given by IDEs and checkers. Additionally, to make use of the pivot model, an extra field `Suspect Lines' will be provided by the pivot model to help LLMs focus on lines with a high possibility of containing dead code.

Since we require a model that is not just a classifier, but also a powerful model that locates, explains, and removes dead code, we turn to LLMs to generate human-readable content that facilitates dead code elimination with minimal human effort. To achieve this, we construct a more detailed and complex training set for LLM fine-tuning.
% The basic dialogue template is presented in Figure \ref{fig1}, and the sources of various data fields are shown in Table \ref{tab1}. Most fields are already populated during the construction of our AIDCE dataset.
% We then leverage the state-of-the-art GPT-4o model as an annotator to complete the unfilled content, including explanations and fixed code snippets.
To be more specific, different from directly querying GPT-4, we provide it with accurate signals including the classification labels and location of dead code, ensuring the output quality of explanations and fixed code.
% Additionally, to make the most of the pivot model, we include an extra field, "Suspect Lines," supplied by the pivot model, to help the LLMs focus on lines with a high likelihood of containing dead code.



\subsubsection{Model Training} We train a pivot model to function as a quick filter and dead code locator, and an LLM as a robust and effective model for double-checking suspect code snippets. The LLM is also trained to generate detailed explanations for dead code and provide practical patches for its elimination.

\textbf{Pivot Model. }
We expect the pivot model to be slim but effective, with a high recall rate in classifying dead code. This ensures that it successfully detects nearly all dead code snippets. While false positives (normal code misclassified as dead code) may occur, the LLMs deployed afterward will perform a double-check mechanism to minimize this. In practice, we choose CodeBERT~\cite{feng2020codebert}, an encoder-only transformer-based pre-trained model, as the backbone.

\textbf{LLMs. }
% Existing powerful LLMs such as GPT-4o can achieve the most part of our task. However, not trained with extensive samples of dead code, even the state-of-the-art LLMs can still ... Moreover, larger LLMs requires more computational resource and results in more API costs. To tackle the issues above, we select Llama-3-8B-Instruct~\cite{} as the base model and make use of our curated synthetic data to perform a supervised fine-tuning step. Note that most fields of our synthetic data are gold labels from the AIDCE dataset, the correctness can be proved, while with the silver labels generated from GPT-4o can enhance the patch generation quality. To achieve this, we employ the LLaMA-Factory~\cite{zheng2024llamafactory} library to train a LoRA adapter with the learning rate of 5e-5 for 3 epochs.
Existing powerful LLMs, such as GPT-4o, can accomplish most of our tasks. However, even state-of-the-art LLMs may struggle with dead code detection due to a lack of extensive training samples specifically for dead code. Additionally, larger LLMs require more computational resources and incur higher API costs. To address these issues, we select Qwen2-7B-Instruct as our base model and utilize our curated synthetic data for supervised fine-tuning (SFT). Our synthetic data is comprised mostly of gold labels from the AIDCE dataset, ensuring accuracy, while silver labels generated by GPT-4o enhance patch generation quality. We employ the LLaMA-Factory~\cite{zheng2024llamafactory} library to train a LoRA~\cite{hulora} adapter.


\subsection{Inference Phase}

% \subsubsection{Dead Code Shapley} Most previous work that leverages language models on code analysis per line such as static program slicing~\cite{yadavally2024learning} and fault localization~\cite{yang2024large} train a line-level encoder to represent each line as an embedding vector. Then a classifier is deployed to make a prediction for every single line. The prediction step only receive the embedding of a single line as input, thus, context in the code snippets are ignored. Especially for our dead code elimination task, 全局层面的信息才是判断某行代码是否包含死代码的关键。比如对于unused代码, 需要关注变量和method在下文中的引用情况，而对于unreachable代码，需要参考上文中对condition中所涉及的关键变量的的赋值和修改。 为了不失去模型对整段代码片段的学习和关注能力，我们将pivot模型设计为一个三分类器，接受整段代码作为输入而仅输出是否包含dead code.

% Most previous work leveraging language models for code analysis on a per-line basis, such as static program slicing~\cite{yadavally2024learning} and fault localization~\cite{yang2024large}, trains a line-level encoder to represent each line as an embedding vector. Then, a classifier is deployed to make predictions for each line individually. This approach only takes the embedding of a single line as input, ignoring the broader context of the code snippets. However, for our dead code elimination task, global-level context is crucial in determining whether a particular line contains dead code. For example, detecting unused code requires tracking references to variables and methods throughout the code, while identifying unreachable code necessitates examining how key variables in conditions are assigned and modified in preceding lines. Thus, to maintain the model's ability to learn from and focus on the entire code snippet, we designed the pivot model as a three-class classifier that accepts the whole code snippet as input and outputs whether it contains dead code. This approach ensures that the model is equipped with a global awareness of the whole code snippet context.

\subsubsection{Dead Code Attribution} Previous work on leveraging language models for per-line code analysis, such as static program slicing~\cite{yadavally2024learning} and fault localization~\cite{yang2024large}, typically uses line-level encoders to represent each line as an embedding vector. These models predict each line individually, ignoring the broader code context. However, dead code elimination requires a global understanding of the code, including tracking variable references and analyzing conditions across multiple lines. To address this, we designed the pivot model as a three-class classifier that processes the entire code snippet, providing a global context for dead code detection but lacking the capability to pinpoint the exact lines.


To empower the pivot model with dead code localization, we introduce a novel method, dead code attribution, which measures the effect of each line on the classification task. Adapting the concept of attribution value (or Shapley value) widely used in model interpretation~\cite{mosca2022shap, nguyen2021effectiveness}, we apply it to code analysis, marking the first use of attribution value in a code-related task to assess the contribution of each line to the model's prediction values instead of single token~\cite{li2016understanding, kim2020interpretation}.



% To be specific, we train a pivot classifier $f$, receiving a $n$-line code snippet $C = {l_1, l_2, \cdots, l_n}$ as input, where $l_i = {c_1^i, c_2^i, \cdots, c_{m_i}^i}$ represents the code tokens of the $i$-th line with total $m_i$ tokens. classifier $f$ makes a prediction among 3 possible classes normal/unused/unreachable to get a probabilities vector $\boldsymbol{a} \in [-1, 1]^3$.

% Before estimating the attribution of each code line $l_i \in C$, we need to examine the 死代码elimination任务中至关重要的特性，that is, 删除死代码对于程序功能是没有影响的，而删除具有实际功能的代码反而可能会产生死代码，因为这同时移除了对相关变量或者method的引用，会导致其他的行成为新的死代码。所以，如果我们对特定的代码行$c_i$执行删除，那么对dead code的prediction的变化可以反映出$l_i$是否包含死代码。如果对应死代码类别的prediction value降低了，那么删除的$l_i$，反之亦然。在实际操作中，我们可以将死代码分为两类。对于非condition类型的代码行，如赋值和计算的相关代码，我们使用Leave-one-out的方法，直接将其删除，即$C{-i} = C - {l_i}.而对于可能被判定为unreachable code的condition类型的代码行，如if(condition)或者while(condition),我们并不能直接删除它，因为这会明显地改变代码的执行逻辑，可能产生新的死代码。对此，我们选择将condition字段替换成CodeBERT模型中的掩码字段[mask],使得可能为恒真或者恒假的condition无法被判断值。这样能够显著降低这段代码被分类为unreachable的可能性。

To be specific, we train a pivot classifier $f$ receiving an $n$-line code snippet $C = \{l_1, l_2, \cdots, l_n\}$ as input, where $l_i = \{c_1^i, c_2^i, \cdots, c_{m_i}^i\}$ represents the tokens of the $i$-th line, containing a total of $m_i$ tokens. The classifier $f$ predicts among 3 possible classes: normal, unused, or unreachable, resulting in a probability vector $f(C) \in [0, 1]^3$.

Before assessing the attribution of each code line $l_i \in C$, it is essential to understand a crucial characteristic in dead code elimination: removing dead code does not affect the program's functionality, whereas removing functional code can create new dead code by eliminating references to variables or methods, causing other lines to become dead code. Thus, if we delete a specific code line $l_i$, the change in the dead code prediction can indicate whether $l_i$ contains dead code. If the prediction value for the dead code class decreases, the deleted $l_i$ was likely dead code, and vice versa. We must point out that, the removal of dead code perfectly fits the motivation for evaluating the importance of each line in classification.

We classify dead code into two types. For non-condition code lines, such as assignment and computation statements, we use the Leave-One-Out (LOO) strategy that removes each line of code:
\begin{equation}
    C_{-i} = C - \{l_i\}
\end{equation}

For condition-type code lines that can be classified into \textbf{unreachable}, like \textit{if(condition)} or \textit{while(condition)}, directly deleting them will significantly alter the code's execution logic and potentially create new dead code. Instead, we replace the conditions with a mask token from the CodeBERT model, \textit{[mask]}, making it impossible to determine the value of conditions as always true or false:
\begin{equation}
    C_{-i} = mask(C, l_i)
\end{equation}
By masking the condition, if the original condition was unreachable, masking it will prevent it from being classified as unreachable. If the condition is normal and functional, masking it will obscure its true behavior and increase the probability of misclassification as unused or unreachable. However, we only need to focus on the \textbf{decrease} of probabilities on unused and unreachable classes as mentioned above.

Thus, we define the attribution value as:
\begin{equation}
    a_{i} = max(f(C) - f(C_{-i}),\ 0)
\end{equation}

% which denotes the probability difference before and after eliminating the assumptive dead code. It is clear to see that the probability difference stands for the attribution of being classified into dead code.
% We sorely focus on the second and third values of $a_{i}$ representing the attribution value of unused and unreachable code respectively. So $a_{i}$ is regarded as a 2-dim vector $a_{i} \in [0, 1]^2$. The greater $a_{i}$ is, the more possibility that line $c_i$ contains dead code. In this way, we successfully distribute the prediction score
% $f(C)$ along all lines of code. This enables us to precisely locate the lines of dead code without specifically training a line-level classifier. Figure \ref{fig4} presents an example of our dead code attribution algorithm.
which denotes the probability difference before and after eliminating the assumed dead code line. It is clear that this value reflects the attribution of the line about being classified as dead code.


\begin{figure}[ht]
\centering
\includegraphics[width=7.35cm]{icse_fig_shapley_value.pdf}
\caption{An illustrative Python method with dead code.}
\label{fig4}
\end{figure}


We focus on the attribution value of unused and unreachable code. Hence, $a_{i}$ is regarded as a 2-dimensional vector $a_{i} \in [0, 1]^2$. The greater the value of $a_{i}$, the higher the likelihood that line $c_i$ contains dead code. This approach allows us to distribute the prediction score $f(C)$ across all lines of code in a mini-batch, enabling us to locate dead code lines without specifically training a line-level classifier. Figure \ref{fig4} presents an example of our dead code attribution algorithm.





% \subsubsection{Pipeline Overview}
% After enabling locating dead code with our proposed dead code attribution approach, we can complete the classification, location, exploration, and patching pipeline for dead code elimination. Concretely, our approach involves three key steps as illustrated in Figure \ref{fig3}:

% \textbf{Pivot Model Filtering. } We start with a series of tokens representing the code snippet $C = \{l_1, l_2, \cdots, l_n\}$. We first query it in our pivot model $f$, to get a prediction among three categories: normal/unused/unreachable. We filter out samples classified as normal to reduce the computational cost for upcoming steps. During training, we ensure the model has an extremely high recall rate for unused and unreachable classes. Thus, almost all potential dead code can proceed to the following steps.

% \textbf{Dead code Locating. }
% In this step, we handle each snippet $C$ marked as a potential dead code. We leverage the dead code attribution approach for each line of code in a $n$-size mini-batch, yielding output attribution scores $\boldsymbol{a} = \{a_1, a_2, \cdots, a_n \}$. For probability differences of unused (the first dimension) and unreachable (the second dimension), we separately 以它们为关键值从大到小排序。举个例子，在处理unused的时候，我们依据第一维对a进行排序，得到了一个ranked list. Despite from 直接返回top-$k$ results, 我们选用了一个伸缩系数$\tau$, 以 $\frac{max \boldsymbol{a^1}}{\tau}$ 为阈值，选出与最大值相近的一系列值，作为可能存在死代码的lines的候选集合。通过改变$\tau$的大小，我们可以在保证仅抽取出少量代码行的同时，得到一个相对高的Dead code Locating的recall值。 在实验中，我们选择了$\tau =5$ 的setting.

% \textbf{Explaining & Eliminating. }
% 通过前面两步的努力，我们为LLMs提供了更丰富的关于dead code的提示，尤其是对dead code行的定位，从而帮助它更精确地找到并且理解code snippets中的dead code, 从而更加关注于Explaining和eliminating dead code的过程。具体来说，我们将死代码候选集作为源代码的辅助信息输入给经过fine-tuning的LLM，并且要求它输出在Figure \ref{fig2}中展示的模板中所要求的信息。 需要指出的是，我们提供的explaination比IDEs和Checkers所提供的信息更详细，更容易阅读。经过这一步，开发者们可以通过explaination信息判断检测出的死代码是否正确，从而接着review所生成的修正后的代码。在LLM所提供信息的帮助下，我们的DCE-LLM pipeline能够最大限度地降低dead code eliminatiopn中的Human effort.

\subsubsection{Pipeline Overview}
By enabling the precise location of dead code with our proposed dead code attribution approach, we can complete the pipeline for dead code elimination, encompassing classification, location, exploration, and patching. Specifically, our approach involves three key steps, as illustrated in Figure \ref{fig3}:

\textbf{Pivot Model Filtering.} We initiate the process with a series of tokens representing the code snippet $C = {l_1, l_2, \cdots, l_n}$. This snippet is input into our pivot model $f$, which predicts one of three categories: normal, unused, or unreachable. To reduce computational overhead for subsequent steps, we filter out samples classified as normal. During the training phase, we ensure the model achieves a high recall rate for the unused and unreachable classes, thereby ensuring that nearly all potential dead code advances to the following stages.

\textbf{Dead Code Locating.}
In this phase, we process each snippet $C$ identified as containing potential dead code. Utilizing our dead code attribution approach, we evaluate each line of code in an $n$-size mini-batch, resulting in attribution scores $\boldsymbol{a} = {a_1, a_2, \cdots, a_n}$. We focus on the probability differences for unused (the first dimension) and unreachable (the second dimension) lines, ranking them separately in descending order. For instance, when addressing unused code, we sort the first dimension of $a$ to produce a ranked list. Instead of merely returning the top-$k$ results, we employ a scaling factor $\tau$, using $\frac{\max \boldsymbol{a^1}}{\tau}$ as a soft threshold to filter out values closer to the maximum, thereby forming a smaller but more accurate candidate set of lines likely to contain dead code.

\textbf{Explaining \& Eliminating.}
Through the preceding steps, we provide large language models (LLMs) with enriched prompts regarding dead code, particularly in identifying specific dead code lines. This enables the LLMs to more accurately identify and understand dead code within code snippets, thereby enhancing the explanation and elimination process. Specifically, we input the candidate set of suspect lines as auxiliary information into a fine-tuned LLM, instructing it to output detailed information . Notably, our explanations are more comprehensive and user-friendly compared to those offered by traditional IDEs and checkers. Following this step, developers can utilize the explanation information to verify the accuracy of the detected dead code and review the generated corrections. Leveraging the detailed insights provided by the LLM, our DCE-LLM pipeline significantly reduces human effort in the dead code elimination process.


\section{Evaluation}
% exp: context length, cross language, fine-tuning
% In this section, we present a comprehensive evaluation of DCE-LLM to answer the following research questions.

% \noindent \textbf{RQ1. How does DCE-LLM compare against prior tools and LLMs?} We evaluate our DCE-LLM approach as a classifier, to classify unused code and unreachable from the same test dataset.

% \noindent \textbf{RQ2. How well can DCE-LLM generate explanations and eliminate dead code?} We evaluate our DCE-LLM approach as a language model, to assess its capability to provide detailed and convincing explanations and practical code patches.

% \noindent \textbf{RQ3. How does each component contribute to DCE-LLM's performance?} We conduct an ablation study to analyze the impact of major components on its performance.

% \noindent \textbf{RQ4. How well does DCE-LLM generalize to other programming languages?} We assess the performance of DCE-LLM on
% several programming languages that are not seen during training, to test the generalization ability of our proposed method.

\subsection{Evaluation Setup}
% \textbf{Dataset.} As there is no prior dataset dedicated to dead code detection, we evaluate the performance of DCE-LLM against various baselines in the test split of our proposed AIDCE dataset. Note that we choose the report from IDEs as the gold label of unused code, the accuracy of IDEs is 100\%, of course. We check how close DCE-LLM performs to IDEs for unused code classification in this setting. We split the test set into Python and Java, with unused code and unreachable code fold, respectively.

% \textbf{Baselines.} We consider two types of baselines. IDEs and checkers provide gold labels for unused code, thus we solely evaluate them in the unreachable setting. For Java, we select IntelliJ IDEA Java IDE and JAVA JIT, the Java just-in-time compiler that can remove dead code in execution time. Considering Python is an interpreted programming language with dynamic semantics, we adopt IntelliJ PyCharm IDE and Vulture as baselines. Another type of baseline is LLMs. We can apply the base prompt template on LLMs to generate answers with the same format of DCE-LLM. We consider several leading models including calling API from \textit{GPT-4o} and \textit{GPT-3.5-turbo}. We also utilize open-source LLMs such as \textit{Llama3-8b-instruct} and \textit{Qwen2-7B-instruct}. We only evaluate 7B-size LLMs due to our limited computing resources.

\textbf{Dataset.} We evaluate the performance of DCE-LLM against various baselines using the test split of our proposed AIDCE dataset. Note that, for unused code, we use the reports from IDEs as the gold standard labels, assuming their accuracy to be 100\%. The test set is divided into Python and Java subsets, each further split into unused code and unreachable code categories.

\noindent \textbf{Baselines.} Considering that DCE-LLM is testing-free and light-weight which does not require any program running time, we only focus on static analyzing methods or compilers. IDEs and checkers serve as the gold standard for unused code, so we only evaluate them in the context of unreachable code detection. For Java, we select IntelliJ IDEA Java IDE and the Java Just-In-Time (JIT) compiler, which can eliminate dead code during execution. Considering Python is an interpreted programming language with dynamic semantics, we adopt PyCharm IDE and Vulture as baselines.

The second type of baseline involves using Large Language Models (LLMs).
% We use the prompt template with LLMs to generate responses formatted similarly to DCE-LLM\footnote{Our task-specific prompt template performs better than standard CoT prompt `Let's think step by step'.}.
Our evaluation includes several leading models, such as API calls to \textit{GPT-4o} and \textit{GPT-3.5-turbo}, as well as open-source LLMs like \textit{Llama3-8b-instruct} and \textit{Qwen2-7B-instruct}. Code-specific LLMs such as \textit{Deepseek-coder-6.7B} and \textit{Mamba-Codestral-7B} are also evaluated.
Due to our limited computing resources, we only assess LLMs up to a maximum size of 8B parameters. For GPT-4o, we also evaluate it with few-shot prompt.
% Considering the length of the source code can heavily increase the input context, which can distract the LLMs as several prior works point out~\cite{xieexplanation, deng2024large}, we conduct evaluation under the zero-shot setting, following other code benchmarks~\cite{chen2021evaluating}.

% \textbf{Metrics.} We adopt the following metrics to access the versatility of DCE-LLM:

% \textit{Classification Metrics.} For RQ1, we measure the classification ability of various methods. Note that IDEs, checkers, and close-source LLMs do not offer probabilistic prediction results, threshold-based metrics such as AUC and ROC Curse do not fit our evaluation. Thus, we adopt standard evaluation metrics for multiclass classification tasks, including the overall metric $Accuracy = \frac{T}{T+F}$, where T and F represent the numbers of correct predictions and wrong predictions. Moreover, we focus on precision $P = \frac{TP}{TP+FP}$, recall $R = \frac{TP}{TP+FN}$, F1-Score $F1 = \frac{2*P*R}{P+R}$ for unused and unreachable classes. Here, TP = True Positives, FP = False Positives, FN = False Negatives, and TN = True Negatives.

% \textit{Generation Metrics.} For RQ2, we measure the quality of generated contexts from LLMs including explanations and repaired code. For explanations, as no valid data source can serve as a proper gold label, we employ three experienced programmers as annotators, to directly compare each output from DCE-LLM, and other LLMs. Annotators select from three options: worse, close, and better (than baseline). For code fixing, we first evaluate if the program keeps the same output before and after removing dead code, to ensure the consistency and correctness of the program. After that, we also perform the contrastive human evaluation process used for evaluating explanations.

% \begin{table*}[h!]
% \centering
% \renewcommand{\arraystretch}{1.15} % Adjust the row height
% \setlength{\tabcolsep}{3.5pt} % Adjust the column separation
% \small
% \caption{Comparison of various baselines (detailed)}
% \begin{tabular}{c  c ccc ccc ccc c}
% \toprule
% \multirow{2}{*}{\textbf{Approach}} & \multirow{2}{*}{\textbf{Language}}  & \multicolumn{3}{c}{\textbf{Unused}} & \multicolumn{3}{c}{\textbf{Unreachable}} & \multicolumn{3}{c}{\textbf{Normal}} & \multirow{2}{*}{\textbf{Accuracy}} \\
% \cmidrule(lr){3-11}

%  &  & \textbf{R} & \textbf{P} & \textbf{F1} & \textbf{R} & \textbf{P} & \textbf{F1} & \textbf{R} & \textbf{P} & \textbf{F1} & \\
% \midrule

% \multirow{3}{*}{IDEs} & Python  & 100.0 & 100.0 & 100.0 & 0 & 100.0 & 0  & 100.0 & 81.40 & 89.73 & 84.28\\
%  & Java  & 100.0 & 100.0 & 100.0 & 23.08 & 100.0 & 35.49  & 100.0 & 91.98 & 95.82 & 94.19\\
%  & \cellcolor{gray!20}Overall  & \cellcolor{gray!20}100.0 & \cellcolor{gray!20}100.0 & \cellcolor{gray!20}100.0 & \cellcolor{gray!20}5.98 & \cellcolor{gray!20}100.0 & \cellcolor{gray!20}9.48 &  \cellcolor{gray!20}100.0 & \cellcolor{gray!20}77.36 & \cellcolor{gray!20}87.24 & \cellcolor{gray!20}87.89\\

% \midrule
% \midrule

% \multirow{3}{*}{Llama-3-8B} &Python   & 43.06 & 15.02 & 22.28 & 34.39 & 29.02 & 31.48 & 31.96 &  73.54 & 44.56 & 32.70\\
%  & Java   & 36.68 & 37.33 & 37.00 & 80.77 & 31.34 & 45.16 & 55.56 & 77.27 & 64.64 & 48.19\\
%  & \cellcolor{gray!20}Overall & \cellcolor{gray!20}39.78 & \cellcolor{gray!20}20.97 & \cellcolor{gray!20}27.46 & \cellcolor{gray!20}44.40 & \cellcolor{gray!20}29.89 & \cellcolor{gray!20}35.73 &\cellcolor{gray!20}40.39 & \cellcolor{gray!20}75.33 & \cellcolor{gray!20}52.58  & \cellcolor{gray!20}38.34 \\

% \midrule

% \multirow{3}{*}{Qwen2-8B} & Python  & 10.65 & 32.39 & 16.03 & 68.25 & 83.23 & 75.00 & 98.18 & 77.68 & 86.74 & 73.88\\
%  & Java   & 6.11 & 77.78 & 11.34 & 69.23 & 73.47 & 71.29 & 98.47 & 70.63 & 82.26 & 67.49\\
%  & \cellcolor{gray!20}Overall   & \cellcolor{gray!20}8.31 & \cellcolor{gray!20}41.57 & \cellcolor{gray!20}13.86 & \cellcolor{gray!20}68.46 & \cellcolor{gray!20}80.88 & \cellcolor{gray!20}74.16  & \cellcolor{gray!20}98.29 & \cellcolor{gray!20}75.00 & \cellcolor{gray!20}85.08 & \cellcolor{gray!20}71.55\\
% \midrule

% \multirow{3}{*}{GPT-3.5}  & Python   & 12.04 & 36.62 & 18.12 & 56.08 & 70.20 & 62.35  & 93.46 & 77.35 & 84.65 & 72.55\\
%  & Java   & 3.49 & 44.44 & 6.48 & 73.08 & 71.70 & 72.38  & 95.64 & 70.69 & 81.30 & 64.44\\
%  & \cellcolor{gray!20}Overall  & \cellcolor{gray!20}7.64 & \cellcolor{gray!20}38.20 & \cellcolor{gray!20}12.73 & \cellcolor{gray!20}59.75 & \cellcolor{gray!20}70.59 & \cellcolor{gray!20}64.72  & \cellcolor{gray!20}94.24 & \cellcolor{gray!20}74.80 & \cellcolor{gray!20}83.40 & \cellcolor{gray!20}69.59\\

% \midrule

% \multirow{3}{*}{GPT-4o} & Python & 68.06 & 33.26 & 44.68 & 84.66 & 57.97 & 68.82 & 68.77 & 91.47 & 78.51 & 62.90 \\
%  & Java  & 27.51 & 45.65 & 34.33 & 76.92 & 37.04 & 50.00   & 74.95 & 74.14 & 74.54 & 57.62\\
%  & \cellcolor{gray!20}Overall  & \cellcolor{gray!20}47.19 & \cellcolor{gray!20}36.21 & \cellcolor{gray!20}40.98 & \cellcolor{gray!20}82.99 & \cellcolor{gray!20}52.08 & \cellcolor{gray!20}64.00   & \cellcolor{gray!20}70.97 & \cellcolor{gray!20}84.06 & \cellcolor{gray!20}76.96 & \cellcolor{gray!20}60.97\\

% \midrule

% \multirow{3}{*}{GPT-4o (3-shot)} & Python & 41.20 & 44.05 & 42.58 & 37.56 & 61.20 & 46.55 & 85.47 & 80.13 & 82.71 & 70.13 \\
%  & Java  & 16.15 & 49.33 & 24.34 & 32.69 & 50.00 & 39.53   & 88.88 & 70.46 & 78.61 & 63.71\\
%  & \cellcolor{gray!20}Overall  & \cellcolor{gray!20}28.31 & \cellcolor{gray!20}45.48 & \cellcolor{gray!20}34.90 & \cellcolor{gray!20}36.51 & \cellcolor{gray!20}58.66 & \cellcolor{gray!20}45.01   & \cellcolor{gray!20}86.69 & \cellcolor{gray!20}76.30 & \cellcolor{gray!20}81.16 & \cellcolor{gray!20}67.79\\

% \midrule
% \midrule

% \multirow{3}{*}{Codellama-7B} &Python   & 97.68 & 18.31 & 30.84 & 55.02 & 19.29 & 28.57 & 0.84 & 63.63 & 1.67 & 10.89\\
%  & Java   & 95.19 & 33.79 & 49.88 & 9.61 & 5.81 & 7.24 & 5.01 & 79.31 & 9.42 & 25.10\\
%  & \cellcolor{gray!20}Overall & \cellcolor{gray!20}96.40 & \cellcolor{gray!20}23.87 & \cellcolor{gray!20}38.26 & \cellcolor{gray!20}45.22 & \cellcolor{gray!20}17.44 & \cellcolor{gray!20}25.17 &\cellcolor{gray!20}2.33 & \cellcolor{gray!20}75.00 & \cellcolor{gray!20}4.52  & \cellcolor{gray!20}16.07 \\

% \midrule

% \multirow{3}{*}{Deepseek-coder-6.7B} &Python   & 0.0 & 0.0 & 0.0 & 4.16 & 16.66 & 6.66 & 96.40 & 72.82 & 82.97 & 70.52\\
%  & Java   & 0.0 & 0.0 & 0.0 & 0.0 & 0.0 & 0.0 & 94.36 & 65.68 & 77.45 & 62.61\\
%  & \cellcolor{gray!20}Overall & \cellcolor{gray!20}0.0 & \cellcolor{gray!20}0.0 & \cellcolor{gray!20}0.0 & \cellcolor{gray!20}2.85 & \cellcolor{gray!20}9.09 & \cellcolor{gray!20}4.34 &\cellcolor{gray!20}95.71 & \cellcolor{gray!20}70.27 & \cellcolor{gray!20}81.04  & \cellcolor{gray!20}67.67 \\

% \midrule

% \multirow{3}{*}{Mamba-Codestral-7B} &Python   & 6.48 & 24.13 & 10.21 & 7.40 & 82.35 & 13.59 & 98.66 & 71.36 & 82.82 & 68.30 \\
%  & Java   & 7.86 & 81.81 & 14.34 & 1.92 & 100.0 & 3.77 & 99.34 & 68.46 & 81.06 & 66.18\\
%  & \cellcolor{gray!20}Overall & \cellcolor{gray!20}7.19 & \cellcolor{gray!20}40.00 & \cellcolor{gray!20}12.19 & \cellcolor{gray!20}6.22 & \cellcolor{gray!20}83.33 & \cellcolor{gray!20}11.58 &\cellcolor{gray!20}98.91 & \cellcolor{gray!20}70.29 & \cellcolor{gray!20}82.18  & \cellcolor{gray!20}67.53 \\

% \midrule

% \multirow{3}{*}{Qwen-2.5-coder-32B} & Python  & 17.59 & 27.74 & 21.53 & 80.95 & 74.63 & 77.66 & 91.89 & 81.61 & 86.45 & 72.05\\
%  & Java  & 10.04 & 79.31 & 17.83 & 78.85 & 75.93 & 77.36 & 96.95 & 71.89 & 82.56 & 67.63\\
%  & \cellcolor{gray!20}Overall  & \cellcolor{gray!20}13.71 & \cellcolor{gray!20}36.75 & \cellcolor{gray!20}19.97 & \cellcolor{gray!20}80.50 & \cellcolor{gray!20}74.90 & \cellcolor{gray!20}77.60 & \cellcolor{gray!20}93.70 & \cellcolor{gray!20}77.28 & \cellcolor{gray!20}84.97 & \cellcolor{gray!20}70.44\\
% \midrule
% \midrule

% \multirow{3}{*}{DCE-LLM} & Python & 98.61 & 89.50 & 93.83 & 97.88 & 97.37 & 97.63 & 99.52 & 99.16 & 99.34 & 97.00 \\
%  & Java & 89.08 & 100.0 & 94.23 & 88.46 & 97.87 & 92.93 & 100.00 & 98.92 & 99.46 & 95.36 \\
%  & \cellcolor{gray!20}Overall & \cellcolor{gray!20}\textbf{93.71} & \cellcolor{gray!20} \textbf{94.34} & \cellcolor{gray!20}\textbf{94.02} & \cellcolor{gray!20}\textbf{95.85} & \cellcolor{gray!20}\textbf{97.47} & \cellcolor{gray!20}\textbf{96.65} & \cellcolor{gray!20}\textbf{99.69} & \cellcolor{gray!20}\textbf{99.07} & \cellcolor{gray!20}\textbf{99.38} & \cellcolor{gray!20}\textbf{96.40} \\


% \bottomrule
% \end{tabular}
% \label{tab3}
% \end{table*}


\label{metric}
\noindent \textbf{Metrics.} We first assess the classification performance of different methods with standard multi-task classification metrics, including accuracy, precision, recall, and F1-Score for the unused and unreachable classes.
Moreover, we assess the quality of generated content with human evaluation, including explanations and repaired code due to the absence of gold standard of generated contents.

% \textbf{Environment. } The entire training and inference phases are conducted using an AMD Ryzen 9 5950X CPU running Ubuntu 23.04 with 128GB RAM and 2 NVIDIA RTX 4090 24GB GPUs. VLLM framework~\cite{kwon2023efficient} is applied for serving our local LLMs. For relatively stable and consistent generation results, we set $temperature = 0.1$. The $max\_token$ value is aligned with the default value of backbone LLMs.

\begin{table*}[h!]
\centering
\renewcommand{\arraystretch}{1} % Adjust the row height
\setlength{\tabcolsep}{5.5pt} % Adjust the column separation
\small
\begin{tabular}{c ccc ccc ccc c}
\toprule
\multirow{2}{*}{\textbf{Approach}} & \multicolumn{3}{c}{\textbf{Unused}} & \multicolumn{3}{c}{\textbf{Unreachable}} & \multicolumn{3}{c}{\textbf{Normal}} & \multirow{2}{*}{\textbf{Accuracy}} \\
\cmidrule(lr){2-10}
 & \textbf{R} & \textbf{P} & \textbf{F1} & \textbf{R} & \textbf{P} & \textbf{F1} & \textbf{R} & \textbf{P} & \textbf{F1} & \\
\midrule

\multirow{1}{*}{IDEs} & 100.0 & 100.0 & 100.0 & 5.98 & 100.0 & 9.48 & 100.0 & 77.36 & 87.24 & 87.89\\

\midrule
\midrule
\multirow{1}{*}{Llama-3-8B} & 39.78 & 20.97 & 27.46 & 44.40 & 29.89 & 35.73 & 40.39 & 75.33 & 52.58 & 38.34 \\

\midrule
\multirow{1}{*}{Qwen2-8B} & 8.31 & 41.57 & 13.86 & 68.46 & 80.88 & 74.16 & 98.29 & 75.00 & 85.08 & 71.55\\

\midrule
\multirow{1}{*}{GPT-3.5} & 7.64 & 38.20 & 12.73 & 59.75 & 70.59 & 64.72 & 94.24 & 74.80 & 83.40 & 69.59\\

\midrule
\multirow{1}{*}{GPT-4o} & 47.19 & 36.21 & 40.98 & 82.99 & 52.08 & 64.00 & 70.97 & 84.06 & 76.96 & 60.97\\

\midrule
\multirow{1}{*}{GPT-4o (3-shot)} & 28.31 & 45.48 & 34.90 & 36.51 & 58.66 & 45.01 & 86.69 & 76.30 & 81.16 & 67.79\\

\midrule
\midrule
\multirow{1}{*}{Codellama-7B} & 96.40 & 23.87 & 38.26 & 45.22 & 17.44 & 25.17 & 2.33 & 75.00 & 4.52 & 16.07 \\

\midrule
\multirow{1}{*}{Deepseek-coder-6.7B} & 0.0 & 0.0 & 0.0 & 2.85 & 9.09 & 4.34 & 95.71 & 70.27 & 81.04 & 67.67 \\

\midrule
\multirow{1}{*}{Mamba-Codestral-7B} & 7.19 & 40.00 & 12.19 & 6.22 & 83.33 & 11.58 & 98.91 & 70.29 & 82.18 & 67.53 \\

\midrule
\multirow{1}{*}{Qwen-2.5-coder-32B} & 13.71 & 36.75 & 19.97 & 80.50 & 74.90 & 77.60 & 93.70 & 77.28 & 84.97 & 70.44\\

\midrule
\midrule
\rowcolor{gray!20} \multirow{1}{*}{DCE-LLM (Ours)} & \cellcolor{gray!20}\textbf{93.71} & \cellcolor{gray!20}\textbf{94.34} & \cellcolor{gray!20}\textbf{94.02} & \cellcolor{gray!20}\textbf{95.85} & \cellcolor{gray!20}\textbf{97.47} & \cellcolor{gray!20}\textbf{96.65} & \cellcolor{gray!20}\textbf{99.69} & \cellcolor{gray!20}\textbf{99.07} & \cellcolor{gray!20}\textbf{99.38} & \cellcolor{gray!20}\textbf{96.40} \\

\bottomrule
\end{tabular}
\caption{Comparison of various baselines}
\label{tab3}
\end{table*}


% \begin{table*}[h!]
% \centering
% \renewcommand{\arraystretch}{1} % Adjust the row height
% \setlength{\tabcolsep}{2.8pt} % Adjust the column separation
% \small
% \caption{Comparison of various baselines}
% \begin{tabular}{c  c ccc ccc ccc c}
% \toprule
% \multirow{2}{*}{\textbf{Approach}} & \multirow{2}{*}{\textbf{Language}}  & \multicolumn{3}{c}{\textbf{Unused}} & \multicolumn{3}{c}{\textbf{Unreachable}} & \multicolumn{3}{c}{\textbf{Normal}} & \multirow{2}{*}{\textbf{Accuracy}} \\
% \cmidrule(lr){3-11}

%  &  & \textbf{R} & \textbf{P} & \textbf{F1} & \textbf{R} & \textbf{P} & \textbf{F1} & \textbf{R} & \textbf{P} & \textbf{F1} & \\
% \midrule

% \multirow{3}{*}{IDEs} & Python  & 100.0 & 100.0 & 100.0 & 0 & 100.0 & 0  & 100.0 & 81.40 & 89.73 & 84.28\\
%  & Java  & 100.0 & 100.0 & 100.0 & 23.08 & 100.0 & 35.49  & 100.0 & 91.98 & 95.82 & 94.19\\
%  & \cellcolor{gray!20}Overall  & \cellcolor{gray!20}100.0 & \cellcolor{gray!20}100.0 & \cellcolor{gray!20}100.0 & \cellcolor{gray!20}5.98 & \cellcolor{gray!20}100.0 & \cellcolor{gray!20}9.48 &  \cellcolor{gray!20}100.0 & \cellcolor{gray!20}77.36 & \cellcolor{gray!20}87.24 & \cellcolor{gray!20}87.89\\

% \midrule
% \midrule

% \multirow{3}{*}{Llama-3-8B} &Python   & 43.06 & 15.02 & 22.28 & 34.39 & 29.02 & 31.48 & 31.96 &  73.54 & 44.56 & 32.70\\
%  & Java   & 36.68 & 37.33 & 37.00 & 80.77 & 31.34 & 45.16 & 55.56 & 77.27 & 64.64 & 48.19\\
%  & \cellcolor{gray!20}Overall & \cellcolor{gray!20}39.78 & \cellcolor{gray!20}20.97 & \cellcolor{gray!20}27.46 & \cellcolor{gray!20}44.40 & \cellcolor{gray!20}29.89 & \cellcolor{gray!20}35.73 &\cellcolor{gray!20}40.39 & \cellcolor{gray!20}75.33 & \cellcolor{gray!20}52.58  & \cellcolor{gray!20}38.34 \\

% \midrule

% \multirow{3}{*}{Qwen2-8B} & Python  & 10.65 & 32.39 & 16.03 & 68.25 & 83.23 & 75.00 & 98.18 & 77.68 & 86.74 & 73.88\\
%  & Java   & 6.11 & 77.78 & 11.34 & 69.23 & 73.47 & 71.29 & 98.47 & 70.63 & 82.26 & 67.49\\
%  & \cellcolor{gray!20}Overall   & \cellcolor{gray!20}8.31 & \cellcolor{gray!20}41.57 & \cellcolor{gray!20}13.86 & \cellcolor{gray!20}68.46 & \cellcolor{gray!20}80.88 & \cellcolor{gray!20}74.16  & \cellcolor{gray!20}98.29 & \cellcolor{gray!20}75.00 & \cellcolor{gray!20}85.08 & \cellcolor{gray!20}71.55\\
% \midrule

% \multirow{3}{*}{GPT-3.5}  & Python   & 12.04 & 36.62 & 18.12 & 56.08 & 70.20 & 62.35  & 93.46 & 77.35 & 84.65 & 72.55\\
%  & Java   & 3.49 & 44.44 & 6.48 & 73.08 & 71.70 & 72.38  & 95.64 & 70.69 & 81.30 & 64.44\\
%  & \cellcolor{gray!20}Overall  & \cellcolor{gray!20}7.64 & \cellcolor{gray!20}38.20 & \cellcolor{gray!20}12.73 & \cellcolor{gray!20}59.75 & \cellcolor{gray!20}70.59 & \cellcolor{gray!20}64.72  & \cellcolor{gray!20}94.24 & \cellcolor{gray!20}74.80 & \cellcolor{gray!20}83.40 & \cellcolor{gray!20}69.59\\

% \midrule

% \multirow{3}{*}{GPT-4o} & Python & 68.06 & 33.26 & 44.68 & 84.66 & 57.97 & 68.82 & 68.77 & 91.47 & 78.51 & 62.90 \\
%  & Java  & 27.51 & 45.65 & 34.33 & 76.92 & 37.04 & 50.00   & 74.95 & 74.14 & 74.54 & 57.62\\
%  & \cellcolor{gray!20}Overall  & \cellcolor{gray!20}47.19 & \cellcolor{gray!20}36.21 & \cellcolor{gray!20}40.98 & \cellcolor{gray!20}82.99 & \cellcolor{gray!20}52.08 & \cellcolor{gray!20}64.00   & \cellcolor{gray!20}70.97 & \cellcolor{gray!20}84.06 & \cellcolor{gray!20}76.96 & \cellcolor{gray!20}60.97\\

% \midrule

% \multirow{3}{*}{GPT-4o (3-shot)} & Python & 41.20 & 44.05 & 42.58 & 37.56 & 61.20 & 46.55 & 85.47 & 80.13 & 82.71 & 70.13 \\
%  & Java  & 16.15 & 49.33 & 24.34 & 32.69 & 50.00 & 39.53   & 88.88 & 70.46 & 78.61 & 63.71\\
%  & \cellcolor{gray!20}Overall  & \cellcolor{gray!20}28.31 & \cellcolor{gray!20}45.48 & \cellcolor{gray!20}34.90 & \cellcolor{gray!20}36.51 & \cellcolor{gray!20}58.66 & \cellcolor{gray!20}45.01   & \cellcolor{gray!20}86.69 & \cellcolor{gray!20}76.30 & \cellcolor{gray!20}81.16 & \cellcolor{gray!20}67.79\\

% \midrule
% \midrule

% \multirow{3}{*}{Codellama-7B} &Python   & 97.68 & 18.31 & 30.84 & 55.02 & 19.29 & 28.57 & 0.84 & 63.63 & 1.67 & 10.89\\
%  & Java   & 95.19 & 33.79 & 49.88 & 9.61 & 5.81 & 7.24 & 5.01 & 79.31 & 9.42 & 25.10\\
%  & \cellcolor{gray!20}Overall & \cellcolor{gray!20}96.40 & \cellcolor{gray!20}23.87 & \cellcolor{gray!20}38.26 & \cellcolor{gray!20}45.22 & \cellcolor{gray!20}17.44 & \cellcolor{gray!20}25.17 &\cellcolor{gray!20}2.33 & \cellcolor{gray!20}75.00 & \cellcolor{gray!20}4.52  & \cellcolor{gray!20}16.07 \\

% \midrule

% \multirow{3}{*}{Deepseek-coder-6.7B} &Python   & 0.0 & 0.0 & 0.0 & 4.16 & 16.66 & 6.66 & 96.40 & 72.82 & 82.97 & 70.52\\
%  & Java   & 0.0 & 0.0 & 0.0 & 0.0 & 0.0 & 0.0 & 94.36 & 65.68 & 77.45 & 62.61\\
%  & \cellcolor{gray!20}Overall & \cellcolor{gray!20}0.0 & \cellcolor{gray!20}0.0 & \cellcolor{gray!20}0.0 & \cellcolor{gray!20}2.85 & \cellcolor{gray!20}9.09 & \cellcolor{gray!20}4.34 &\cellcolor{gray!20}95.71 & \cellcolor{gray!20}70.27 & \cellcolor{gray!20}81.04  & \cellcolor{gray!20}67.67 \\

% \midrule

% \multirow{3}{*}{Mamba-Codestral-7B} &Python   & 6.48 & 24.13 & 10.21 & 7.40 & 82.35 & 13.59 & 98.66 & 71.36 & 82.82 & 68.30 \\
%  & Java   & 7.86 & 81.81 & 14.34 & 1.92 & 100.0 & 3.77 & 99.34 & 68.46 & 81.06 & 66.18\\
%  & \cellcolor{gray!20}Overall & \cellcolor{gray!20}7.19 & \cellcolor{gray!20}40.00 & \cellcolor{gray!20}12.19 & \cellcolor{gray!20}6.22 & \cellcolor{gray!20}83.33 & \cellcolor{gray!20}11.58 &\cellcolor{gray!20}98.91 & \cellcolor{gray!20}70.29 & \cellcolor{gray!20}82.18  & \cellcolor{gray!20}67.53 \\

% \midrule
% \midrule

% \multirow{3}{*}{DCE-LLM} & Python & 98.61 & 89.50 & 93.83 & 97.88 & 97.37 & 97.63 & 99.52 & 99.16 & 99.34 & 97.00 \\
%  & Java & 89.08 & 100.0 & 94.23 & 88.46 & 97.87 & 92.93 & 100.00 & 98.92 & 99.46 & 95.36 \\
%  & \cellcolor{gray!20}Overall & \cellcolor{gray!20}\textbf{93.71} & \cellcolor{gray!20} \textbf{94.34} & \cellcolor{gray!20}\textbf{94.02} & \cellcolor{gray!20}\textbf{95.85} & \cellcolor{gray!20}\textbf{97.47} & \cellcolor{gray!20}\textbf{96.65} & \cellcolor{gray!20}\textbf{99.69} & \cellcolor{gray!20}\textbf{99.07} & \cellcolor{gray!20}\textbf{99.38} & \cellcolor{gray!20}\textbf{96.40} \\


% \bottomrule
% \end{tabular}
% \label{tab3}
% \end{table*}



\subsection{Classification Performance}
\label{unreach}
% Despite the powerfulness showed in unused code detection task, IDEs, checkers and compilers, they show poor performance when facing unreachable adversarial attacks. Besides the patterns design in DaK~\cite{gao2023discrete}, we extend the number of patterns to 62. Table \ref{tab2} presents several of those patterns. \textit{PyCharm} and \textit{IntelliJ IDEA} serve as IDEs for checking unreachable code. The brown {\textcolor{brown}{\ding{51}}}/{\textcolor{brown}{\ding{55}}}) denotes if they successfully catch the unreachable code in Python and Java, respectively. Meanwhile, we also adopt \textit{Vulture} checker (for Python) and \textit{JIT} compiler (for Java) to check the same pattern set. Their results are provided with the blue {\textcolor{darkblue}{\ding{51}}}/{\textcolor{darkblue}{\ding{55}}}). Note that even the modern and one of the most powerful IDE among our 测试的 tools, IntelliJ IDEA, can only handle 18 of the 62 patterns, which is less than 30\%.


Despite their effectiveness in detecting unused code, IDEs, checkers, and compilers perform poorly when faced with unreachable adversarial attacks. In addition to the patterns designed in DaK~\cite{gao2023discrete}, we have expanded the number of patterns to 62. Table \ref{taba} presents several of these patterns. We use \textit{PyCharm} and \textit{IntelliJ IDEA} as IDEs for checking unreachable code, denoted by brown {\textcolor{brown}{\ding{51}}}/{\textcolor{brown}{\ding{55}}} symbols for Python and Java, respectively. Additionally, we employ the \textit{Vulture} checker (for Python) and the \textit{JIT} compiler (for Java) to evaluate the same pattern set, with their results indicated by blue {\textcolor{darkblue}{\ding{51}}}/{\textcolor{darkblue}{\ding{55}}} symbols. Remarkably, even IntelliJ IDEA, one of the most advanced and widely used IDEs, successfully identifies only 18 out of the 62 patterns, which is less than 30\%. Other static analysis tools such as UCDetector and J2ObjC show weaker detection capabilities compared to IntelliJ IDEA. The UCDetector detected 22 out of 100 unused code cases but failed to detect any of the 62 unreachable patterns.
Considering that those tools are based on rules, we can assume that the precision value $P$ for unreachable code detection is $1$, and the recall rate equals the checked ratio as we randomly select inserted patterns.


% As for approaches based on deep learning, Table \ref{tab3} details the experimental results of how our proposed DCE-LLM compares with other LLMs, we combine IDEs for different languages as a powerful ensemble model. Four LLMs also compete with our proposed method. Except for the ~\textit{unused} split where IDEs serve as gold labels, we can observe that DCE-LLM achieves the highest performance among all metrics in both Python and Java code corpus.

% Due to the unbalanced label distribution in our test set, a model can easily achieve high accuracy by predicting all codes as normal. Even GPT-4o gets a relatively lower accuracy compared with GPT-3.5. Thus, focusing on detecting dead code, the recall and F1 score for unused and unreachable samples deserves more concentration. All LLMs struggle to detect unused codes with recall rates lower than 50\%. The global 特质 of unused code detecting challanges LLMs‘ long-form memorisition ablity, which is still an crucial issue under research. DCE-LLM achieves a aaa\% recall rate, introducing a great improvement compared with other baselines.
% 而对于Unrechable detection来说, LLMs take advantages in analylizing and calculating representations in the condition. While IDEs can barely find unreachable conditions, GPT-4o can extract unreachable code with 82.99\ recall rate, and Qwen2-8B acquires an 80.88\% precision. Our proposed DCE-LLM strongly outperforms them, with a aaa\% leading in recall and a bbb\% 领先 in precision.
For learning-based approaches, Table \ref{tab3} details the experimental results comparing DCE-LLM with other LLMs. IDEs for different languages are ensembled into a powerful baseline. Four foundation LLMs and 3 code-specific LLMs are evaluated alongside our method. Except for the \textit{unused} split where IDEs serve as gold labels, we observe that DCE-LLM achieves the highest scores across all metrics in Python and Java code corpus
\footnote{Detailed performance in Python/Java split are provided in Appendix.}
, including foundation LLMs and code-specific LLMs.

Due to the unbalanced label distribution, a model can easily achieve high accuracy by predicting all code as normal. Even GPT-4o exhibits relatively lower accuracy compared to GPT-3.5. Thus, focusing on the detection of dead code, the recall and F1 scores for unused and unreachable samples are more critical metrics. All LLMs struggle to detect unused code, with recall rates below 50\%. The global characteristics of unused code detection challenge the long-form memorization ability of LLMs, which remains an active area of research. DCE-LLM, however, achieves a recall rate of \textit{93.71\%}, marking a substantial improvement over other baselines.

For unreachable code detection, LLMs demonstrate an advantage in analyzing and computing representations in conditional statements. While IDEs rarely detect unreachable conditions, GPT-4o identifies unreachable code with a recall rate of 82.99\%, and Qwen2-8B achieves a precision of 80.88\%.
Our proposed DCE-LLM surpasses these models, with a leading recall rate of 95.85\% and a precision rate of 97.47\%.


From our study, the most challenging pattern for Python involves type-based unreachable conditions (e.g., \texttt{isInstance(int(a))}). Additionally, the model occasionally misses several dead code lines when multiple dead code lines appear together.




% \begin{custombox}
% As a classifier, DCE-LLM performs near IDEs in unused code detection and performs all baselines in unreachable code detection.
% \end{custombox}

\begin{figure}[ht]
\centering
\includegraphics[width=7.5cm]{icse_fig_win_rate.pdf}
\caption{Human evaluation on explanation generation.}
\label{fig5}
\end{figure}


\begin{figure}[htbp]
\centering
\includegraphics[width=7.5cm]{icse_fig_code.pdf}
\caption{Human evaluation on code repairing.}
\label{fig6}
\end{figure}

\subsection{Generation Quality}

To measure the quality of generated contents from DCE-LLM, we conduct a human evaluation of both explanations and repaired code. Three experienced programmers serve as voluntary annotators on 20 unused codes and 20 unreachable codes.
% We engage three experienced programmers as annotators to directly compare outputs from DCE-LLM and other LLMs. For explainations, annotators rate each output as worse, similar, or better than the baseline. For code fixing, we first ask annotators to check the consistency of execution before and after dead code removal to ensure correctness. Then,
% For explanation generating, we select 20 unused codes and 20 unreachable codes in Python, then evaluate the quality of the explanation by directly comparing them with the IDE standard description such as \textit{`Unused import statement from sys import exit'}, and GPT-4o. Annotators can select better/close/worse than baseline.

% For generated codes, we further judge if the generated code is equivalent to the original one as dead code does not affect the execution result. As depicted in Figure \ref{fig5}, our method outperformed the baseline methods in explanation. IDEs obtain precision while lacking diversity in explanation and cannot handle most unreachable statements, while DCE-LLM aligns well with the annotation from GPT-4o.

For explanation, annotators rate each output as worse, similar, or better than the baseline. As depicted in Figure \ref{fig5}, our method outperformed the baseline methods in explanation.

For code repairing, we focus on two major points: (1) if the code maintains the same function and (2) if all dead code is removed. Only samples satisfying both two conditions are positive. As shown in Figure \ref{fig6}, DCE-LLM surpasses both Qwen and GPT-4o in the evaluation.

Our study revealed that IDEs provide only template-based explanations for unused code, whereas our model generates more diverse explanations. Furthermore, due to the high-quality synthetic data and attribution techniques used, our model can accurately locate and explain dead code lines, leading to more applicable code patches.


% To be more specific,
% GPT-4o is more aggressive in removing key expressions such as \textit{stdin = sys.stdin}, while ours does not. The removal of this kind of expression can result in different functions of the program. Some dead code remains when handling long and complex programs.


% \begin{custombox}
% As a generation model, DCE-LLM outputs high-quality explanations and repaired code to reduce the human efforts in dead code elimination.
% \end{custombox}



\subsection{Ablation Study}
To understand the impact of the major components in DCE-LLM, we separate our pivot model and generation model to test them as individual classifiers. Additionally, we explore minor modifications of the pipeline. Our ablation results are presented in Table \ref{tab4}, demonstrating the effectiveness of the essentials in our pipeline.

\begin{table}[h]
\centering
\small
\renewcommand{\arraystretch}{1.2}
\resizebox{\columnwidth}{!}{
\begin{tabular}{l ccc ccc}

\toprule
\multirow{2}{*}{\textbf{Approach}} & \multicolumn{3}{c}{\textbf{Unused}} & \multicolumn{3}{c}{\textbf{Unreachable}} \\
\cmidrule(lr){2-7}
 & \textbf{R} & \textbf{P} & \textbf{F1} & \textbf{R} & \textbf{P} & \textbf{F1} \\

\midrule

\multirow{1}{*}{\raggedright \cellcolor{gray!20}\ \ DCE-LLM} & \cellcolor{gray!20}93.71 & \cellcolor{gray!20}94.34 & \cellcolor{gray!20}94.02 & \cellcolor{gray!20}95.85 & \cellcolor{gray!20}97.47 & \cellcolor{gray!20}96.65 \\

\multirow{1}{*}{\raggedright - Pivot Model} & 51.46 & 87.07 & 64.69 & 84.65 & 99.02 & 91.28 \\
\multirow{1}{*}{\raggedright - LLM} & 94.61 & 92.32 & 93.45 & 94.19 & 99.13 & 96.60 \\
\multirow{1}{*}{\raggedright - SFT} & 48.52 & 96.57 & 64.59 & 98.58 & 50.73 & 66.99 \\
\multirow{1}{*}{\raggedright - Attribution} & 83.37 & 98.40 & 90.27 & 91.30 & 95.85 & 93.52 \\
\bottomrule
\end{tabular}}
\caption{Ablation study results of DCE-LLM}
\label{tab4}
\end{table}


% \textit{Without Pivot Model, }we directly prompt our fine-tuned LLM to conduct dead code elimination. We observe great delegation (about 30\% for F1 score) in unused code detection and a slight drop for unreachable code as well. Our pivot model can retrieve most dead code snippets from the corpus.

% \textit{Without LLM, }the classification results changes slightly. However, the capability of generating explanations and repairing code is lost.
\textit{Without Pivot Model, }we directly prompt our fine-tuned LLM to perform dead code elimination. This results in a significant decline in performance, with a drop of about 30\% in the F1 score for unused code detection and a slight decrease for unreachable code. The pivot model plays a crucial role in retrieving most dead code snippets from the corpus.

\textit{Without LLM, }the classification results show only minor changes. However, the ability to generate explanations and repair code is completely lost.

\textit{Without SFT, }we replace our fine-tuned LLM with Qwen2-7B-instruct model. This substitution results in a 30\% drop in the F1 score for both splits. Without the specialized training on our well-annotated corpus, the LLM cannot neither maintaining high classification performance or generate high-quality contents.

\textit{Without Dead Code Attribution, }the LLM loses precise guidance on where dead code is located, leading to confusion in the dead code elimination process. In this configuration, all metrics decrease by approximately 3\%.

% \begin{custombox}
% The pivot model contributes to the classification result. The LLM contributes to generation capability while preserving outstanding classification performance.
% \end{custombox}




\subsection{Programming Language Generalization}
To test the cross-language generalization of DCE-LLM, we annotated 275 Golang files as a test set. We applied a unreachable pattern attack on this set, resulting in 100 unused and 98 unreachable samples (some of which overlap).
% Note that this process is different from the dataset construction for Java/Python,
The Golang dataset is simpler than Java/Python with less random generated variables and obfuscations.
We then leveraged DCE-LLM for dead code detection on this corpus, and the results are reported in Table \ref{tab5}. From the table, we can observe that even when facing an unseen programming language, DCE-LLM still showcases strong performance in detecting unused code.
For unreachable code, the performance remains robust, although slightly lower compared to unused code detection. This indicates that while DCE-LLM effectively generalizes to new programming languages, there is still room for improvement in handling more complex unreachable code scenarios. Performance on both unused and unreachable code outperforms GPT-4o.


\begin{table}[h]
\centering
\small
\renewcommand{\arraystretch}{1.2}
\resizebox{\columnwidth}{!}{
\begin{tabular}{l ccc ccc}

\toprule
\multirow{2}{*}{\textbf{Approach}} & \multicolumn{3}{c}{\textbf{Unused}} & \multicolumn{3}{c}{\textbf{Unreachable}} \\
\cmidrule(lr){2-7}
 & \textbf{R} & \textbf{P} & \textbf{F1} & \textbf{R} & \textbf{P} & \textbf{F1} \\

\midrule

\multirow{1}{*}{GPT-4o} & 73.00 & 62.93 & 67.59 & 80.61 & 83.16 & 81.87 \\
\multirow{1}{*}{\cellcolor{gray!20}DCE-LLM} & \cellcolor{gray!20}85.00 & \cellcolor{gray!20}96.59 & \cellcolor{gray!20}90.42 & \cellcolor{gray!20}90.00 & \cellcolor{gray!20}91.84 & \cellcolor{gray!20}90.91 \\

\bottomrule
\end{tabular}}
\caption{Performance of DCE-LLM on Golang}
\label{tab5}
\end{table}

% \begin{table}[ht]
% \small
% \renewcommand{\arraystretch}{1}
%     \caption{Performance of DCE-LLM on Golang}
%     \centering

%     \begin{tabular}{cccc}
%     \toprule
%            & R & P & F1 \\
%     \midrule
%         Unused & 85.00 & 96.59 & 90.42\\
%         Unreachable & 90.00& 91.84& 90.91\\
%         Normal & 87.25&  100.0& 93.19\\
%         \bottomrule
%     \end{tabular}
%     \label{tab5}
% \end{table}
% \begin{table}[ht]
% \small
% \renewcommand{\arraystretch}{1}
%     \caption{Performance comparison of DCE-LLM and DCE-LLM2 on Golang}
%     \centering

%     \begin{tabular}{ccccc}
%     \toprule
%            & \multicolumn{2}{c}{DCE-LLM} & \multicolumn{2}{c}{DCE-LLM2} \\
%            \cmidrule(lr){2-3} \cmidrule(lr){4-5}
%            & R & F1 & R & F1 \\
%     \midrule
%         Unused      & 85.00 & 90.42 & 88.00 & 91.00 \\
%         Unreachable & 90.00 & 90.91 & 91.00 & 92.50 \\
%         Normal      & 87.25 & 93.19 & 89.50 & 94.00 \\
%     \bottomrule
%     \end{tabular}
%     \label{tab5}
% \end{table}


\section{Conclusion}

In this paper, we proposed DCE-LLM, the first learning-based solution for the dead code elimination task that leverages a pivot model for filtering and recall, and an LLM for explanation and code fixing. The LLM is empowered with location information provided by the dead code attribution technique.
By fine-tuning both the pivot model and LLM on the first large-scale dead code dataset, we successfully constructed a complete pipeline for dead code elimination, encompassing classification, location, explanation, and repair. Our proposed framework demonstrates superior performance in both classification and generation tasks. Additionally, the versatility of DCE-LLM enables generalization on other programming languages.



% For future work, we aim to support longer contexts and complex control flows for method-level and repository-level dead code elimination and further enhance the robustness and efficiency of our framework. We also plan to combine it with program analysis techniques to enhance DCE-LLM in diverse development environments.

% \begin{custombox}
% On unseen programming languages, DCE-LLM demonstrates sufficient versatility to effectively perform DCE.
% \end{custombox}

% \subsection{Discussion}
% During our experiments, we observed that even the gold labels from IDEs are sometimes incorrect. A notable example is the handling of unused variables in the \textit{\_\_main\_\_} namespace in Python, such as \textit{a = 1, b = 2, print(a)}. In this case, the dead code line \textit{b = 2} is not reported by PyCharm unless it is defined within a method, making it a local variable. In contrast, our DCE-LLM can easily identify such dead code. Additionally, DCE-LLM supports code snippets with minor bugs that prevent them from being compiled, further showcasing its robustness and versatility.

% line location

% is ide right?

% \section{Discussion}



\section*{Limitation}
Firstly, the pivot model introduces a relatively limited input length. The input length is 512 for CodeBERT, while programs can easily exceed. Thus, this work focus on line-level dead code elimination instead of method-level. We plan to support longer contexts and complex control flows in the future.

Secondly, the study of prompting techniques can be insufficient. We only compares the standard CoT prompt with ours. We did not studied prompt engineering and devising more sophisticated in-context examples which is beyond the scope of this work. We also believe that no prompting method can bridge the over 30\% F1 performance gap.

Lastly, as a learning-based method, DCE-LLM is not a complete that detects all dead code. We plan to conduct case study on failure cases to improve our model in the future.


\bibliography{custom}

\newpage
\appendix

\section{Dataset Details}
The AIDCE dataset is divided into training, validation, and test subsets with 15,043, 1,889, and 1,891 samples, respectively. To accurately represent the dead code rate from the original CodeNet dataset while including unreachable codes, we set the ratio of normal to unused to unreachable code at around 4:1:1. Some overlapped samples (around 5\%) that contain both unused and unreachable code also exist.

\begin{table}[ht]
    \centering
    \begin{tabular}{>{\centering\arraybackslash}m{0.8cm} >{\centering\arraybackslash}m{2cm} >{\centering\arraybackslash}m{2.8cm}}
        \toprule
        \multicolumn{2}{c}{\textbf{Fields}} & \textbf{Source} \\
        \midrule
        \multirow{2}{*}{Input} & Source code & Dataset \\
                               & Suspect lines & Pivot Model\\
        \cmidrule(lr){1-3}
        \multirow{5}{*}{Output} & Dead code & Dataset \\
                                & Line number & Dataset \\
                                & Type & Dataset \\
                                & Explanation & Dataset \& GPT-4o \\
                                & Fixed code & Source Code \& GPT-4o \\
        \bottomrule
    \end{tabular}
    \caption{Source of fields in SFT data for training LLMs}
    \label{taba1}
\end{table}

For SFT dataset, We collect 1,500 samples with a ratio of normal to unused to unreachable code at about 1:1:1 as our pivot model can filter out most normal code snippets before passing them into LLMs. Table \ref{taba1} provides details of source of the fields for SFT data.


\section{Model Training}
 For CodeBERT model, we a classification head on top of CodeBERT. The model is fine-tuned to predict the class of code snippets among three categories. We set the batch size to 16 and the learning rate to 5e-5. The model is fine-tuned using the Adam optimizer for up to 3 epochs, with an early stopping mechanism to prevent overfitting.

 For LLM training, we employ the LLaMA-Factory~\cite{zheng2024llamafactory} WEB GUI, we fine-tuned a LoRA~\cite{hulora} for the Qwen2-7B model with the following settings: learning rate of 5e-5, 10 epochs, a cosine learning rate scheduler, LoRA rank of 8, and alpha of 16. Additionally, we used DeepSpeed ZeRO-2 for efficient GPU memory management.

\section{Experimental Environment}
The entire training and inference phases are conducted using an AMD Ryzen 9 5950X CPU running Ubuntu 23.04 with 128GB RAM and 2 NVIDIA RTX 4090 24GB GPUs. VLLM framework~\cite{kwon2023efficient} is applied for serving our local LLMs. For relatively stable and consistent generation results, we set $temperature = 0.1$. The $max\_token$ value is aligned with the default value of backbone LLMs. In our experiments, adjusting $\tau$ allows us to extract a small number of code lines while maintaining a high recall rate for dead code location. We set the soft threshold parameter $\tau = 2$.

\section{Detailed Performance}
We present the detailed performance in Tab. \ref{taba3}.


\begin{table*}[h!]
\centering
\renewcommand{\arraystretch}{1.15} % Adjust the row height
\setlength{\tabcolsep}{3.5pt} % Adjust the column separation
\small
\caption{Comparison of various baselines (detailed)}
\begin{tabular}{c  c ccc ccc ccc c}
\toprule
\multirow{2}{*}{\textbf{Approach}} & \multirow{2}{*}{\textbf{Language}}  & \multicolumn{3}{c}{\textbf{Unused}} & \multicolumn{3}{c}{\textbf{Unreachable}} & \multicolumn{3}{c}{\textbf{Normal}} & \multirow{2}{*}{\textbf{Accuracy}} \\
\cmidrule(lr){3-11}

 &  & \textbf{R} & \textbf{P} & \textbf{F1} & \textbf{R} & \textbf{P} & \textbf{F1} & \textbf{R} & \textbf{P} & \textbf{F1} & \\
\midrule

\multirow{3}{*}{IDEs} & Python  & 100.0 & 100.0 & 100.0 & 0 & 100.0 & 0  & 100.0 & 81.40 & 89.73 & 84.28\\
 & Java  & 100.0 & 100.0 & 100.0 & 23.08 & 100.0 & 35.49  & 100.0 & 91.98 & 95.82 & 94.19\\
 & \cellcolor{gray!20}Overall  & \cellcolor{gray!20}100.0 & \cellcolor{gray!20}100.0 & \cellcolor{gray!20}100.0 & \cellcolor{gray!20}5.98 & \cellcolor{gray!20}100.0 & \cellcolor{gray!20}9.48 &  \cellcolor{gray!20}100.0 & \cellcolor{gray!20}77.36 & \cellcolor{gray!20}87.24 & \cellcolor{gray!20}87.89\\

\midrule
\midrule

\multirow{3}{*}{Llama-3-8B} &Python   & 43.06 & 15.02 & 22.28 & 34.39 & 29.02 & 31.48 & 31.96 &  73.54 & 44.56 & 32.70\\
 & Java   & 36.68 & 37.33 & 37.00 & 80.77 & 31.34 & 45.16 & 55.56 & 77.27 & 64.64 & 48.19\\
 & \cellcolor{gray!20}Overall & \cellcolor{gray!20}39.78 & \cellcolor{gray!20}20.97 & \cellcolor{gray!20}27.46 & \cellcolor{gray!20}44.40 & \cellcolor{gray!20}29.89 & \cellcolor{gray!20}35.73 &\cellcolor{gray!20}40.39 & \cellcolor{gray!20}75.33 & \cellcolor{gray!20}52.58  & \cellcolor{gray!20}38.34 \\

\midrule

\multirow{3}{*}{Qwen2-8B} & Python  & 10.65 & 32.39 & 16.03 & 68.25 & 83.23 & 75.00 & 98.18 & 77.68 & 86.74 & 73.88\\
 & Java   & 6.11 & 77.78 & 11.34 & 69.23 & 73.47 & 71.29 & 98.47 & 70.63 & 82.26 & 67.49\\
 & \cellcolor{gray!20}Overall   & \cellcolor{gray!20}8.31 & \cellcolor{gray!20}41.57 & \cellcolor{gray!20}13.86 & \cellcolor{gray!20}68.46 & \cellcolor{gray!20}80.88 & \cellcolor{gray!20}74.16  & \cellcolor{gray!20}98.29 & \cellcolor{gray!20}75.00 & \cellcolor{gray!20}85.08 & \cellcolor{gray!20}71.55\\
\midrule

\multirow{3}{*}{GPT-3.5}  & Python   & 12.04 & 36.62 & 18.12 & 56.08 & 70.20 & 62.35  & 93.46 & 77.35 & 84.65 & 72.55\\
 & Java   & 3.49 & 44.44 & 6.48 & 73.08 & 71.70 & 72.38  & 95.64 & 70.69 & 81.30 & 64.44\\
 & \cellcolor{gray!20}Overall  & \cellcolor{gray!20}7.64 & \cellcolor{gray!20}38.20 & \cellcolor{gray!20}12.73 & \cellcolor{gray!20}59.75 & \cellcolor{gray!20}70.59 & \cellcolor{gray!20}64.72  & \cellcolor{gray!20}94.24 & \cellcolor{gray!20}74.80 & \cellcolor{gray!20}83.40 & \cellcolor{gray!20}69.59\\

\midrule

\multirow{3}{*}{GPT-4o} & Python & 68.06 & 33.26 & 44.68 & 84.66 & 57.97 & 68.82 & 68.77 & 91.47 & 78.51 & 62.90 \\
 & Java  & 27.51 & 45.65 & 34.33 & 76.92 & 37.04 & 50.00   & 74.95 & 74.14 & 74.54 & 57.62\\
 & \cellcolor{gray!20}Overall  & \cellcolor{gray!20}47.19 & \cellcolor{gray!20}36.21 & \cellcolor{gray!20}40.98 & \cellcolor{gray!20}82.99 & \cellcolor{gray!20}52.08 & \cellcolor{gray!20}64.00   & \cellcolor{gray!20}70.97 & \cellcolor{gray!20}84.06 & \cellcolor{gray!20}76.96 & \cellcolor{gray!20}60.97\\

\midrule

\multirow{3}{*}{GPT-4o (3-shot)} & Python & 41.20 & 44.05 & 42.58 & 37.56 & 61.20 & 46.55 & 85.47 & 80.13 & 82.71 & 70.13 \\
 & Java  & 16.15 & 49.33 & 24.34 & 32.69 & 50.00 & 39.53   & 88.88 & 70.46 & 78.61 & 63.71\\
 & \cellcolor{gray!20}Overall  & \cellcolor{gray!20}28.31 & \cellcolor{gray!20}45.48 & \cellcolor{gray!20}34.90 & \cellcolor{gray!20}36.51 & \cellcolor{gray!20}58.66 & \cellcolor{gray!20}45.01   & \cellcolor{gray!20}86.69 & \cellcolor{gray!20}76.30 & \cellcolor{gray!20}81.16 & \cellcolor{gray!20}67.79\\

\midrule
\midrule

\multirow{3}{*}{Codellama-7B} &Python   & 97.68 & 18.31 & 30.84 & 55.02 & 19.29 & 28.57 & 0.84 & 63.63 & 1.67 & 10.89\\
 & Java   & 95.19 & 33.79 & 49.88 & 9.61 & 5.81 & 7.24 & 5.01 & 79.31 & 9.42 & 25.10\\
 & \cellcolor{gray!20}Overall & \cellcolor{gray!20}96.40 & \cellcolor{gray!20}23.87 & \cellcolor{gray!20}38.26 & \cellcolor{gray!20}45.22 & \cellcolor{gray!20}17.44 & \cellcolor{gray!20}25.17 &\cellcolor{gray!20}2.33 & \cellcolor{gray!20}75.00 & \cellcolor{gray!20}4.52  & \cellcolor{gray!20}16.07 \\

\midrule

\multirow{3}{*}{Deepseek-coder-6.7B} &Python   & 0.0 & 0.0 & 0.0 & 4.16 & 16.66 & 6.66 & 96.40 & 72.82 & 82.97 & 70.52\\
 & Java   & 0.0 & 0.0 & 0.0 & 0.0 & 0.0 & 0.0 & 94.36 & 65.68 & 77.45 & 62.61\\
 & \cellcolor{gray!20}Overall & \cellcolor{gray!20}0.0 & \cellcolor{gray!20}0.0 & \cellcolor{gray!20}0.0 & \cellcolor{gray!20}2.85 & \cellcolor{gray!20}9.09 & \cellcolor{gray!20}4.34 &\cellcolor{gray!20}95.71 & \cellcolor{gray!20}70.27 & \cellcolor{gray!20}81.04  & \cellcolor{gray!20}67.67 \\

\midrule

\multirow{3}{*}{Mamba-Codestral-7B} &Python   & 6.48 & 24.13 & 10.21 & 7.40 & 82.35 & 13.59 & 98.66 & 71.36 & 82.82 & 68.30 \\
 & Java   & 7.86 & 81.81 & 14.34 & 1.92 & 100.0 & 3.77 & 99.34 & 68.46 & 81.06 & 66.18\\
 & \cellcolor{gray!20}Overall & \cellcolor{gray!20}7.19 & \cellcolor{gray!20}40.00 & \cellcolor{gray!20}12.19 & \cellcolor{gray!20}6.22 & \cellcolor{gray!20}83.33 & \cellcolor{gray!20}11.58 &\cellcolor{gray!20}98.91 & \cellcolor{gray!20}70.29 & \cellcolor{gray!20}82.18  & \cellcolor{gray!20}67.53 \\

\midrule

\multirow{3}{*}{Qwen-2.5-coder-32B} & Python  & 17.59 & 27.74 & 21.53 & 80.95 & 74.63 & 77.66 & 91.89 & 81.61 & 86.45 & 72.05\\
 & Java  & 10.04 & 79.31 & 17.83 & 78.85 & 75.93 & 77.36 & 96.95 & 71.89 & 82.56 & 67.63\\
 & \cellcolor{gray!20}Overall  & \cellcolor{gray!20}13.71 & \cellcolor{gray!20}36.75 & \cellcolor{gray!20}19.97 & \cellcolor{gray!20}80.50 & \cellcolor{gray!20}74.90 & \cellcolor{gray!20}77.60 & \cellcolor{gray!20}93.70 & \cellcolor{gray!20}77.28 & \cellcolor{gray!20}84.97 & \cellcolor{gray!20}70.44\\
\midrule
\midrule

\multirow{3}{*}{DCE-LLM} & Python & 98.61 & 89.50 & 93.83 & 97.88 & 97.37 & 97.63 & 99.52 & 99.16 & 99.34 & 97.00 \\
 & Java & 89.08 & 100.0 & 94.23 & 88.46 & 97.87 & 92.93 & 100.00 & 98.92 & 99.46 & 95.36 \\
 & \cellcolor{gray!20}Overall & \cellcolor{gray!20}\textbf{93.71} & \cellcolor{gray!20} \textbf{94.34} & \cellcolor{gray!20}\textbf{94.02} & \cellcolor{gray!20}\textbf{95.85} & \cellcolor{gray!20}\textbf{97.47} & \cellcolor{gray!20}\textbf{96.65} & \cellcolor{gray!20}\textbf{99.69} & \cellcolor{gray!20}\textbf{99.07} & \cellcolor{gray!20}\textbf{99.38} & \cellcolor{gray!20}\textbf{96.40} \\


\bottomrule
\end{tabular}
\label{taba3}
\end{table*}


\section{Example Prompts}

\subsection{Base Prompt}

\begin{figure}[ht]
\begin{minipage}{\linewidth}
\small
\noindent\makebox[\linewidth]{\rule{\linewidth}{1pt}}
\colorbox{codebg}{
\parbox{\dimexpr\linewidth-2\fboxsep}{
\textbf{\textcolor{cyan}{Input:}} \\
You are an experienced programmer expert in detecting dead code such as unused code and unreachable code. Given the following code, respond in the following format: \\
\textcolor{orange}{Dead code:} Yes or No \\
For each dead code,\\
\textcolor{gray}{Line Number:} Line number\\
\textcolor{olive}{Type:} Unused or Unreachable\\
\textcolor{purple}{Explanation:} Explanation for prediction\\

Finally, fix the code with: \\
\textcolor{brown}{Fixed Code:} Fixed code \\

\textcolor{blue}{Code:} [Code in Figure 1] \\

\textbf{\textcolor{magenta}{Output:}} \\
\textcolor{orange}{Dead code:} \textcolor{red}{Yes} \\
\textcolor{gray}{Line Number:} \textcolor{red}{4} \\
\textcolor{olive}{Type:} Unused \\
\textcolor{purple}{Explanation:} The variable s3 is defined but never used in any subsequent code. Instead of using s3, the code mistakenly uses the literal string `s3'. \\

\textcolor{brown}{Fixed Code:} [Code in Figure 1 after fix line 11]
}}
\noindent\makebox[\linewidth]{\rule{\linewidth}{1.5pt}}
\end{minipage}
\end{figure}


\subsection{Prompt with Suspect Lines}
\begin{figure}[ht]
\begin{minipage}{\linewidth}
\small
\noindent\makebox[\linewidth]{\rule{\linewidth}{1pt}}
\colorbox{codebg}{
\parbox{\dimexpr\linewidth-2\fboxsep}{
\textbf{\textcolor{cyan}{Input:}} \\
You are an experienced programmer expert in detecting dead code such as unused code and unreachable code. Given the following code and \textcolor{red}{suspect lines}, respond in the following format: \\
\textcolor{orange}{Dead code:} Yes or No \\
For each dead code,\\
\textcolor{gray}{Line Number:} Line number\\
\textcolor{olive}{Type:} Unused or Unreachable\\
\textcolor{purple}{Explanation:} Explanation for prediction\\

Finally, fix the code with: \\
\textcolor{brown}{Fixed Code:} Fixed code \\

\textcolor{blue}{Code:} [Code in Figure 1] \\
\textcolor{darkblue}{Suspect Lines:} [Suspect lines from the pivot model] \\


\textbf{\textcolor{magenta}{Output:}} \\
\textcolor{orange}{Dead code:} \textcolor{red}{Yes} \\
\textcolor{gray}{Line Number:} \textcolor{red}{4} \\
\textcolor{olive}{Type:} Unused \\
\textcolor{purple}{Explanation:} The variable s3 is defined but never used in any subsequent code. Instead of using s3, the code mistakenly uses the literal string `s3'. \\

\textcolor{brown}{Fixed Code:} [Code in Figure 1 after fix line 11]
}}
\noindent\makebox[\linewidth]{\rule{\linewidth}{1.5pt}}
\end{minipage}
\end{figure}



\end{document}
